% ==========================================================
% 第14章 代数幾何とゲージ理論
% ==========================================================
\chapter{代数幾何とゲージ理論}
\label{ch:ch14}

\begin{chapterlead}
ゲージ理論の幾何学的な実体は，ファイバー束とその上の接続である．ゲージ場は接続，場の強さは曲率，
物質場は同伴ベクトル束の切断，ゲージ変換は束の自己同型であり，物理の言葉と幾何の言葉はほぼ一対一に対応する．
本章ではまず主束・ベクトル束・接続・曲率というゲージ理論の幾何学的基礎を整え，
次に底空間が複素多様体の場合に，\textbf{正則ベクトル束}という代数幾何的な対象と
\textbf{接続}という微分幾何的な対象がどう結びつくかを調べる．
その帰結が「曲率の $(0,2)$ 成分が消える接続 ＝ 正則構造」という対応と，
「Hermitian--Yang--Mills 方程式の解 ＝ 安定な正則ベクトル束」という Donaldson--Uhlenbeck--Yau の定理である．
さらに Higgs 場を加えた Higgs 束と Hitchin 方程式を導入し，ゲージ同値類のなす空間（モジュライ空間）へと話をつなげる．
Chern 類や Kähler 幾何の詳細は第15章，モジュライ空間の一般論は第16章で扱う．
\end{chapterlead}

% ----------------------------------------------------------
\section{ゲージ理論の幾何学的基礎}
% ----------------------------------------------------------

\subsection{ゲージ理論の幾何学的定式化}

電磁気学から始めよう．荷電粒子の波動関数 $\psi(x)$ は位相の付け替え $\psi\mapsto e^{i\theta(x)}\psi$ に対して物理的に不変であり，
微分 $\partial_\mu\psi$ はこの変換で共変でないため，ゲージ場 $A^{\rm phys}_\mu$ を導入して共変微分
$D_\mu = \partial_\mu - i g A^{\rm phys}_\mu$ を用いる．ゲージ変換は
\begin{equation}
  \psi\mapsto e^{i\theta}\psi,\qquad A^{\rm phys}_\mu\mapsto A^{\rm phys}_\mu + \frac1g\partial_\mu\theta
\end{equation}
である．幾何学的に読み直すと次のようになる：
\begin{itemize}
  \item 各点 $x$ で $\psi(x)$ は「位相の基準を決めて初めて複素数になる」量，すなわち時空の各点に付随した1次元複素ベクトル空間 $E_x$ の元である．$E_x$ を束ねたものが\textbf{ベクトル束} $E\to X$，$\psi$ はその\textbf{切断}である．
  \item 位相の基準の選び方（局所自明化）は点ごとに自由であり，その自由度 $U(1)$ が\textbf{構造群}である．
  \item 異なる点のファイバー $E_x,E_{x+dx}$ を比較する規則が\textbf{接続}であり，$A_\mu$ はその局所表示である．
  \item 場の強さ $F_{\mu\nu}$ は無限小ループを一周したときの位相のずれ，すなわち\textbf{曲率}である．
\end{itemize}

\textbf{記法の約束}：本章では数学の流儀にならい，接続形式 $A$ を Lie 代数 $\mathfrak{g}$ に値をとる1形式とし，共変微分を
\begin{equation}
  D = d + A,\qquad A = -i g A^{\rm phys}_\mu dx^\mu
  \label{eq:ch14-convention}
\end{equation}
と書く．$U(n)$ ゲージ理論では $\mathfrak{g}=\mathfrak{u}(n)$ は反 Hermite 行列からなり，物理で用いる Hermite 行列値の $A^{\rm phys}$ と
$-ig$ 倍だけ異なる．曲率も $F = dA + A\wedge A = -ig F^{\rm phys}$ となる（$F^{\rm phys}_{\mu\nu}=\partial_\mu A^{\rm phys}_\nu-\partial_\nu A^{\rm phys}_\mu - ig[A^{\rm phys}_\mu,A^{\rm phys}_\nu]$）．

\subsection{主束}

\begin{mydef}{主束}{ch14-principal}
$G$ を Lie 群とする．\textbf{主 $G$ 束} $\pi:P\to X$ とは，多様体 $P$ と滑らかな全射 $\pi$，
および $P$ への右作用 $P\times G\to P$，$(p,g)\mapsto p\cdot g$ の組で，次をみたすものをいう．
\begin{enumerate}
  \item $G$ の作用はファイバーを保ち（$\pi(p\cdot g)=\pi(p)$），各ファイバー $P_x=\pi^{-1}(x)$ 上で自由かつ推移的である．
  \item （局所自明性）$X$ の開被覆 $\{U_\alpha\}$ と $G$ 同変な微分同相 $\phi_\alpha:\pi^{-1}(U_\alpha)\to U_\alpha\times G$ が存在する．ここで $U_\alpha\times G$ には $G$ が右から $(x,h)\cdot g=(x,hg)$ と作用する．
\end{enumerate}
$X$ を底空間，$P$ を全空間，$G$ を構造群という．
\end{mydef}

各ファイバー $P_x$ は $G$ と同相であるが，「単位元」に相当する特別な点は指定されていない
（$G$-torsor）．$P_x$ の点を一つ選ぶことは，ゲージ（基準となる枠）を一つ選ぶことに相当する．

\begin{figure}[ht]
\centering
\begin{tikzpicture}[scale=1.0]
  % base
  \draw[thick] (0,0) .. controls (2,-0.4) and (4,0.4) .. (6,0) node[right]{$X$};
  % fibers
  \foreach \x/\y in {0.8/-0.12,2/-0.13,3.2/0.0,4.4/0.13} {
    \draw[blue!70!black,thick] (\x,\y+0.5) -- (\x,\y+2.6);
    \fill (\x,\y) circle (1.2pt);
  }
  \node[blue!70!black] at (4.4,3.0) {$P_x\cong G$};
  \node at (4.4,-0.35) {$x$};
  \draw[->,thick] (5.4,1.6) -- (5.4,0.4) node[midway,right]{$\pi$};
  % section
  \draw[red!70!black,thick] (0.8,1.2) .. controls (1.6,1.9) and (2.6,0.9) .. (4.4,1.7);
  \node[red!70!black] at (2.5,1.95) {局所切断 $s_\alpha$};
  \draw[thick,gray!70!black] (0.6,-0.4)--(0.6,-0.5)--(3.4,-0.5)--(3.4,-0.4);
  \node[below] at (2.0,-0.5) {$U_\alpha$};
\end{tikzpicture}
\caption{主束 $P\to X$ の模式図．各点 $x$ の上にファイバー $P_x\cong G$ が立ち，局所切断 $s_\alpha$ はゲージの選択を表す．}
\label{fig:ch14-bundle}
\end{figure}

\begin{myex}{主束の例}{ch14-principal-ex}
(1) 自明束 $X\times G$．
(2) Hopf 束 $S^3\to S^2=\PP^1$：$S^3=\{(Z_0,Z_1):|Z_0|^2+|Z_1|^2=1\}$ に $U(1)$ が $(Z_0,Z_1)e^{i\theta}$ で作用し，商が $\PP^1$．
これは磁荷1の Dirac モノポールの主 $U(1)$ 束である．同様に $\C^{n+1}\setminus\{0\}\to\PP^n$ は主 $\C^*$ 束．
(3) 枠束 $F(X)$：$n$ 次元多様体 $X$ の各点の接空間の基底全体．主 $\GL(n,\R)$ 束である．
(4) 等質空間 $G\to G/H$（$H$ は閉部分群）は主 $H$ 束．
\end{myex}

\subsection{構造群}

構造群 $G$ は「ファイバーの内部自由度を互いに移す変換の群」であり，物理ではゲージ群である．
標準模型では $G=SU(3)\times SU(2)\times U(1)$，一般相対論の局所 Lorentz 対称性では $SO(1,3)$ が構造群になる．

構造群は必要に応じて小さくしたり（\textbf{構造群の簡約}），大きくしたりできる．
たとえば計量を入れたベクトル束では $\GL(n,\C)$ を $U(n)$ に簡約でき，
これが Hermite 計量の選択と同値である（14.4.3節）．逆に $U(n)\subset\GL(n,\C)$ を複素化することは，
「ユニタリなゲージ変換」から「複素ゲージ変換」へ拡大することにあたり，本章後半で重要な役割を果たす．

構造群の Lie 代数 $\mathfrak{g}$ には随伴表現 $\mathrm{Ad}:G\to\GL(\mathfrak{g})$，$\mathrm{Ad}_g\xi = g\xi g^{-1}$ が作用する．
接続形式と曲率は $\mathfrak{g}$ に値をとり，ゲージ変換で随伴的に変換する．
\begin{note}
Lie 代数の約束：$\mathfrak{u}(n)=\{\xi\in M_n(\C):\xi^\dagger=-\xi\}$，$\mathfrak{su}(n)$ はそのうちトレース零のもの．
物理では生成子 $T^a$ を Hermite にとり $\xi = -i\theta^aT^a$ と書く．\eqref{eq:ch14-convention} の $-ig$ はこの違いである．
\end{note}

\subsection{局所自明化}

局所自明化 $\phi_\alpha:\pi^{-1}(U_\alpha)\to U_\alpha\times G$ は，局所切断 $s_\alpha(x)=\phi_\alpha^{-1}(x,e)$ を与えることと同値である．
二つの局所自明化が重なる $U_\alpha\cap U_\beta$ では，$s_\beta(x) = s_\alpha(x)\cdot g_{\alpha\beta}(x)$ をみたす滑らかな写像
\begin{equation}
  g_{\alpha\beta}:U_\alpha\cap U_\beta\to G
\end{equation}
が一意に定まる．これを\textbf{遷移関数}という．定義から直ちに次の\textbf{コサイクル条件}がしたがう：
\begin{equation}
  g_{\alpha\alpha}=e,\qquad g_{\alpha\beta}\,g_{\beta\gamma} = g_{\alpha\gamma}\quad (U_\alpha\cap U_\beta\cap U_\gamma \text{ 上}),
  \qquad g_{\beta\alpha} = g_{\alpha\beta}^{-1}.
  \label{eq:ch14-cocycle}
\end{equation}

\begin{figure}[ht]
\centering
\begin{tikzpicture}[scale=1.0]
  \draw[thick] (0,0) -- (8,0) node[right]{$X$};
  \draw[thick,blue!70!black] (0.3,-0.3) -- (5.0,-0.3);
  \node[blue!70!black,below] at (1.5,-0.3) {$U_\alpha$};
  \draw[thick,red!70!black] (3.0,-0.6) -- (7.7,-0.6);
  \node[red!70!black,below] at (6.5,-0.6) {$U_\beta$};
  \fill[gray!20] (3.0,0.4) rectangle (5.0,3.0);
  \node at (4.0,3.25) {\small $U_\alpha\cap U_\beta$};
  \draw[blue!70!black,thick] (0.3,1.4) .. controls (2,1.2) and (3.5,1.0) .. (5.0,1.3);
  \node[blue!70!black] at (1.2,1.65) {$s_\alpha$};
  \draw[red!70!black,thick] (3.0,2.4) .. controls (4.5,2.6) and (6,2.0) .. (7.7,2.3);
  \node[red!70!black] at (6.8,2.55) {$s_\beta$};
  \draw[-{Stealth},thick] (4.0,1.13) -- (4.0,2.5) node[midway,right]{\small $\cdot g_{\alpha\beta}(x)$};
  \foreach \x in {0.6,1.8,3.0,4.0,5.0,6.2,7.4} \draw[gray] (\x,0.4)--(\x,3.0);
\end{tikzpicture}
\caption{局所自明化と遷移関数．$U_\alpha\cap U_\beta$ 上で2つの局所切断（ゲージの選択）は $s_\beta=s_\alpha g_{\alpha\beta}$ で結ばれる．}
\label{fig:ch14-transition}
\end{figure}

逆に，開被覆と \eqref{eq:ch14-cocycle} をみたす $\{g_{\alpha\beta}\}$ が与えられれば，
$\bigsqcup_\alpha U_\alpha\times G$ を $(x,h)_\beta\sim(x,g_{\alpha\beta}(x)h)_\alpha$ で貼り合わせて主束が再構成できる．
局所切断を $s'_\alpha = s_\alpha h_\alpha$ と取り替えると $g'_{\alpha\beta}=h_\alpha^{-1}g_{\alpha\beta}h_\beta$ となるので，
\begin{center}
主 $G$ 束の同型類 $\longleftrightarrow$ $\{g_{\alpha\beta}\}$ のコサイクル条件をみたす族 $/$ $\{h_\alpha\}$ による変換
\end{center}
であり，これは Čech コホモロジー $H^1(X,\underline{G})$（$\underline{G}$ は $G$ 値滑らかな関数の層）にほかならない．
これは第11章の層コホモロジーの言葉で，$G=\C^*$，正則関数に限れば $H^1(X,\OO_X^*)=\Pic(X)$（正則直線束の同型類）になる．

\begin{myex}{Dirac モノポール}{ch14-dirac}
$S^2$ を北半球 $U_N$ と南半球 $U_S$ で覆うと，重なりは赤道の帯 $\simeq S^1$ である．主 $U(1)$ 束は
遷移関数 $g_{NS}(\varphi)=e^{in\varphi}$（$n\in\Z$）で決まり，$n$ が磁荷（Dirac の量子化条件）である．
$g_{NS}$ が $S^1\to U(1)$ の写像として一価であることが $n\in\Z$ を強制する．
$n=\pm1$ がHopf 束に対応する（符号は向きの約束による）．
\end{myex}

\subsection{ゲージ変換}

\begin{mydef}{ゲージ変換}{ch14-gauge}
主 $G$ 束 $P\to X$ の\textbf{ゲージ変換}とは，$G$ 同変でファイバーを保つ微分同相 $f:P\to P$（$f(p\cdot g)=f(p)\cdot g$，$\pi\circ f=\pi$）をいう．
ゲージ変換全体は群をなし，\textbf{ゲージ群} $\mathcal{G}$ と書く．
\end{mydef}

局所切断 $s_\alpha$ を用いると $f(s_\alpha(x)) = s_\alpha(x)\cdot h_\alpha(x)$ と書け，$h_\alpha:U_\alpha\to G$ は重なりで
$h_\alpha = g_{\alpha\beta}h_\beta g_{\alpha\beta}^{-1}$ をみたす．すなわちゲージ変換は随伴束 $\mathrm{Ad}P=P\times_G G$ の切断である．
自明束では $\mathcal{G}=C^\infty(X,G)$ であり，物理でいう「点ごとに異なる $G$ の変換 $g(x)$」である．

ゲージ群 $\mathcal{G}$ は無限次元の群であるが，大域的な対称性（物理的に観測される対称性）ではなく
\textbf{記述の冗長性}である．物理的な配位空間はゲージ同値類 $\mathcal{A}/\mathcal{G}$（$\mathcal{A}$ は接続全体）であり，
これがモジュライ空間（14.6節）の原型となる．

\begin{physbox}[ゲージ変換は冗長性]
「ゲージ対称性」は対称性ではなく自由度の数え過ぎである，とよく言われる．
幾何学的には，物理的な情報は束 $P$ と接続の同型類にあり，局所自明化（ゲージ）の選択は座標系の選択と同様である．
一方，位相的に非自明な束（モノポール，インスタントン）では大域的な自明化そのものが存在せず，
ゲージを一つの座標で固定しようとすると Dirac 弦のような特異点が現れる．
\end{physbox}

\exprob{ch14-pr-cocycle}
コサイクル条件 \eqref{eq:ch14-cocycle} から $g_{\alpha\alpha}=e$ と $g_{\beta\alpha}=g_{\alpha\beta}^{-1}$ を導け．
また局所切断を $s'_\alpha=s_\alpha h_\alpha$ と取り替えたとき，$g'_{\alpha\beta}=h_\alpha^{-1}g_{\alpha\beta}h_\beta$ もコサイクル条件をみたすことを確かめよ．

% ----------------------------------------------------------
\section{ベクトル束}
% ----------------------------------------------------------

\subsection{ベクトル束の定義}

\begin{mydef}{ベクトル束}{ch14-vb}
階数 $r$ の（実または複素）\textbf{ベクトル束} $\pi:E\to X$ とは，各ファイバー $E_x=\pi^{-1}(x)$ が
$r$ 次元ベクトル空間であり，開被覆 $\{U_\alpha\}$ とファイバーごとに線形な微分同相
$\phi_\alpha:\pi^{-1}(U_\alpha)\to U_\alpha\times\mathbb{K}^r$（$\mathbb{K}=\R$ または $\C$）が存在するものをいう．
階数1のベクトル束を\textbf{直線束}という．
\end{mydef}

主束とベクトル束は次のように行き来できる．主 $G$ 束 $P$ と表現 $\rho:G\to\GL(V)$ に対し，
\begin{equation}
  E = P\times_\rho V := (P\times V)/\bigl((p\cdot g,v)\sim(p,\rho(g)v)\bigr)
\end{equation}
は $V$ をファイバーとするベクトル束になる．これを\textbf{同伴束}という．
遷移関数は $\rho(g_{\alpha\beta})$ である．逆に階数 $r$ のベクトル束 $E$ からは，各ファイバーの基底全体として
主 $\GL(r)$ 束（枠束）$\mathrm{Fr}(E)$ が得られ，$E=\mathrm{Fr}(E)\times_{\GL(r)}\mathbb{K}^r$ である．

物理的には，ゲージ群 $G$ の同じ主束 $P$ に対し，物質場の属する表現ごとに異なる同伴束が得られる．
たとえば $SU(3)$ の基本表現（クォーク）と随伴表現（グルーオン）は同じ $P$ の異なる同伴束の切断である．

\subsection{局所自明化}

ベクトル束の局所自明化は，$U_\alpha$ 上の局所枠（局所的に各点で基底をなす切断の組）$e^{(\alpha)}=(e^{(\alpha)}_1,\dots,e^{(\alpha)}_r)$ を選ぶことと同値である．
重なりでは
\begin{equation}
  e^{(\beta)}_j = \sum_i e^{(\alpha)}_i\,(g_{\alpha\beta})_{ij},\qquad g_{\alpha\beta}:U_\alpha\cap U_\beta\to\GL(r,\mathbb{K})
\end{equation}
となり，切断 $s=\sum_i s_\alpha^i e^{(\alpha)}_i$ の成分は
\begin{equation}
  s_\alpha = g_{\alpha\beta}\,s_\beta
  \label{eq:ch14-section-transform}
\end{equation}
と変換する．遷移関数はやはりコサイクル条件 \eqref{eq:ch14-cocycle} をみたす．

ベクトル束に対する線形代数の操作はすべて遷移関数の操作として実現される：
\begin{center}
\begin{tabular}{c|c}
束 & 遷移関数 \\ \hline
直和 $E\oplus F$ & $g^E_{\alpha\beta}\oplus g^F_{\alpha\beta}$ \\
テンソル積 $E\otimes F$ & $g^E_{\alpha\beta}\otimes g^F_{\alpha\beta}$ \\
双対 $E^*$ & ${}^t(g_{\alpha\beta})^{-1}$ \\
行列式束 $\det E=\bigwedge^r E$ & $\det g_{\alpha\beta}$ \\
自己準同型束 $\End E = E^*\otimes E$ & $\xi\mapsto g_{\alpha\beta}\xi g_{\alpha\beta}^{-1}$
\end{tabular}
\end{center}
特に直線束の遷移関数は $\mathbb{K}^*$ 値関数であり，テンソル積は遷移関数の積に対応する．

\subsection{ベクトル束の射}

\begin{mydef}{ベクトル束の射}{ch14-vb-map}
同じ底空間上のベクトル束 $E,F$ の間の\textbf{射} $\varphi:E\to F$ とは，ファイバーを保ち（$\pi_F\circ\varphi=\pi_E$）
各ファイバー上で線形な滑らかな写像をいう．局所自明化のもとで $\varphi$ は行列値関数 $\varphi_\alpha:U_\alpha\to M_{r_F\times r_E}$ で表され，
$\varphi_\alpha g^E_{\alpha\beta} = g^F_{\alpha\beta}\varphi_\beta$ をみたす．
逆射をもつ射を同型といい，$\varphi$ は $\Hom(E,F)=E^*\otimes F$ の切断と同一視できる．
\end{mydef}

射の核や像は，階数が一定であればまたベクトル束になる．しかし階数が点によって跳ぶと束にならない．
たとえば $X=\C$ 上の自明直線束の射 $\varphi:\underline{\C}\to\underline{\C}$，$v\mapsto zv$ は
$z=0$ でのみ零写像であり，余核は原点に集中した「摩天楼層」$\C_0=\OO/z\OO$ になる．
このような現象を扱うためにベクトル束を\textbf{連接層}（第11章）にまで広げる必要がある．
階数 $r$ のベクトル束は局所自由な $\OO_X$ 加群（局所的に $\OO_X^{\oplus r}$ と同型な層）と同じものである．

\subsection{切断}

\begin{mydef}{切断}{ch14-section}
ベクトル束 $E\to X$ の（滑らかな）\textbf{切断}とは，滑らかな写像 $s:X\to E$ で $\pi\circ s=\mathrm{id}_X$ をみたすものをいう．
切断全体 $\Gamma(X,E)$ はベクトル空間（$C^\infty(X)$ 加群）をなす．局所的には $s$ は成分 $s_\alpha:U_\alpha\to\mathbb{K}^r$ の族であり，
\eqref{eq:ch14-section-transform} で貼り合わさる．
\end{mydef}

ベクトル束の切断は常に存在する（零切断）．一方，主束の切断の存在は束が自明であることと同値である
（$s$ があれば $p=s(x)g\mapsto (x,g)$ が自明化を与える）．また，零点をもたない切断が存在するかどうかは位相的な問題であり，
たとえば $TS^2$ には零点のないベクトル場は存在しない（毛玉定理）．これは Euler 標数 $\chi(S^2)=2\neq0$ の反映である．

\begin{physbox}[物質場は同伴束の切断]
表現 $\rho$ の物質場（スカラー場・Dirac 場）$\phi$ は，局所的には $V$ 値関数 $\phi_\alpha(x)$ であるが，
異なるゲージの間で $\phi_\alpha=\rho(g_{\alpha\beta})\phi_\beta$ と変換する．これはまさに同伴束 $E=P\times_\rho V$ の切断である．
「場は関数である」という素朴な見方が許されるのは束が自明な場合に限られ，
モノポール背景の荷電粒子の波動関数は「切断」であって関数ではない（Wu--Yang の定式化）．
\end{physbox}

\subsection{ベクトル束の例}

\begin{myex}{Möbius 帯}{ch14-mobius}
$S^1$ を2つの弧 $U_1,U_2$ で覆うと，重なりは2つの小区間 $I_+\sqcup I_-$ からなる．遷移関数を
$g_{12}=+1$（$I_+$ 上），$g_{12}=-1$（$I_-$ 上）と定めると，実直線束 $M\to S^1$ が得られる．これが Möbius 帯である．
$M$ の任意の切断は少なくとも1点で零になり（中間値の定理），$M$ は自明束 $S^1\times\R$ と同型でない．
$M\otimes M$ は遷移関数 $(\pm1)^2=1$ で自明となる．
\end{myex}

\begin{figure}[ht]
\centering
\begin{tikzpicture}[scale=1.0]
  \fill[blue!8] (0,0) rectangle (5,1.6);
  \draw[thick] (0,0)--(5,0); \draw[thick] (0,1.6)--(5,1.6);
  \draw[thick,-{Stealth}] (0,0)--(0,1.6);
  \draw[thick,-{Stealth}] (5,1.6)--(5,0);
  \node[left] at (0,0.8) {$A$}; \node[right] at (5,0.8) {$A$};
  \draw[red!70!black,thick] (0,0.4) .. controls (1.8,0.2) and (3.2,1.4) .. (5,1.2);
  \fill[red!70!black] (2.5,0.8) circle (1.5pt);
  \node[red!70!black,above] at (2.5,0.9) {\small 零点};
  \draw[dashed] (0,0.8)--(5,0.8);
  \node at (2.5,-0.35) {\small 底空間 $S^1=[0,1]/0\sim1$};
\end{tikzpicture}
\caption{Möbius 帯．左右の辺を向きを反転させて貼り合わせる．切断（赤）は左端の値 $s$ が右端で $-s$ に対応するので必ず零点（破線との交点）をもつ．}
\label{fig:ch14-mobius}
\end{figure}

\begin{myex}{$S^2=\PP^1$ の接束}{ch14-ts2}
$\PP^1$ を座標 $z$（$U_0$）と $w=1/z$（$U_1$）で覆い，局所枠を $e_0=\partial_z$，$e_1=\partial_w$ とする．
$\partial_w = \frac{dz}{dw}\partial_z = -z^2\,\partial_z$ であるから，$e_1=e_0g_{01}$ の形で正則接束 $T\PP^1$ の遷移関数は
$g_{01}=-z^2$ であり，後で見るように $T\PP^1\cong\OO(2)$ である．ベクトル場 $\partial_z$ は $w=0$（無限遠点）で2位の零点をもち，
零点の総数 $2$ は Euler 標数 $\chi(S^2)=2$ に一致する（Poincaré--Hopf）．
\end{myex}

\begin{myex}{トートロジカル直線束 $\OO(-1)$}{ch14-taut}
$\PP^n$ の点 $\ell=[Z]$ は $\C^{n+1}$ の原点を通る直線 $\ell$ である．ファイバーとして $\ell$ そのものをとった
\begin{equation}
  \OO(-1) = \{(\ell,v)\in\PP^n\times\C^{n+1} : v\in\ell\}
\end{equation}
を\textbf{トートロジカル直線束}という．$U_\alpha=\{Z_\alpha\neq0\}$ 上の局所枠を $e_\alpha([Z])=Z/Z_\alpha\in\ell$ とすれば，
$e_\alpha = (Z_\beta/Z_\alpha)\,e_\beta$ であるから，遷移関数（$e_\beta = e_\alpha g_{\alpha\beta}$）は
\begin{equation}
  g_{\alpha\beta} = \frac{Z_\alpha}{Z_\beta}
\end{equation}
であり，正則である．双対束 $\OO(1)=\OO(-1)^*$ の遷移関数は $Z_\beta/Z_\alpha$，そのテンソル冪 $\OO(d)=\OO(1)^{\otimes d}$ の遷移関数は
$(Z_\beta/Z_\alpha)^d$ である．$\OO(d)$（$d\geq0$）の正則切断は $d$ 次同次多項式 $P(Z)$ と同一視でき（局所的に $P(Z)/Z_\alpha^d$），
$\dim H^0(\PP^n,\OO(d)) = \binom{n+d}{d}$ となる．$\OO(-1)$ を $\C^{n+1}\times\PP^n$ から第1成分に射影すると，
$\C^{n+1}$ の原点での爆発（blow-up）が得られる（第15章・第17章）．
\end{myex}

$\PP^1$ の場合，$\OO(d)$ の遷移関数は $g_{01}=(Z_1/Z_0)^{d}=z^d$（座標 $z=Z_1/Z_0$）であり，
例\ref{ex:ch14-ts2} の $T\PP^1$ の遷移関数 $-z^{2}$ と比べて（定数 $-1$ は枠の取り替えで吸収できる）$T\PP^1\cong\OO(2)$ が分かる．
双対の余接束（＝標準束）は $K_{\PP^1}\cong\OO(-2)$ である．

\exprob{ch14-pr-od}
$\PP^1$ 上で $\OO(d)$ の正則切断が $d$ 次同次多項式 $P(Z_0,Z_1)$ と一対一に対応することを示し，$\dim H^0(\PP^1,\OO(d))=d+1$（$d\geq0$），$H^0(\PP^1,\OO(d))=0$（$d<0$）を確かめよ．

\exprob{ch14-pr-mobius}
Möbius 帯 $M\to S^1$ の任意の切断が零点をもつことを示せ．また $M\oplus M$ は自明束 $S^1\times\R^2$ と同型であることを示せ．
（ヒント：$M\oplus M\cong M\otimes\C$ を複素直線束とみよ．）

% ----------------------------------------------------------
\section{接続と曲率}
% ----------------------------------------------------------

\subsection{接続の幾何学的定義}

ベクトル束 $E$ の異なるファイバーは自然には比較できない．比較の規則を与えるのが接続である．

\begin{mydef}{接続（共変微分）}{ch14-connection}
ベクトル束 $E\to X$ の\textbf{接続}とは，$\mathbb{K}$-線形写像
\begin{equation}
  D:\Gamma(X,E)\to\Gamma(X,T^*X\otimes E)=\Omega^1(X,E)
\end{equation}
で，Leibniz 則 $D(fs) = df\otimes s + f\,Ds$（$f\in C^\infty(X)$）をみたすものをいう．
\end{mydef}

主束の言葉では，接続は $P$ の各点の接空間を「垂直方向」（ファイバー方向）と「水平方向」に $G$ 同変に分解すること，
あるいは同値に $\mathfrak{g}$ 値1形式 $\omega\in\Omega^1(P,\mathfrak{g})$（接続形式）を与えることである．
水平方向は「隣のファイバーへ移るときに内部自由度を変えない方向」を定める．
接続があれば，$X$ 上の曲線 $\gamma$ に沿ってファイバーを移す\textbf{平行移動} $P_\gamma:E_{\gamma(0)}\to E_{\gamma(1)}$ が
常微分方程式 $D_{\dot\gamma}s=0$ の解として定まる．

\subsection{共変微分}

局所枠 $e^{(\alpha)}$ を固定すると，$D e^{(\alpha)}_j = \sum_i e^{(\alpha)}_i\,(A_\alpha)_{ij}$ により行列値1形式 $A_\alpha$ が定まり，
切断の成分に対して
\begin{equation}
  D s_\alpha = d s_\alpha + A_\alpha s_\alpha,\qquad D_\mu = \partial_\mu + A_\mu
\end{equation}
と書ける．\eqref{eq:ch14-convention} により $D_\mu = \partial_\mu - igA^{\rm phys}_\mu$ であり，物理の共変微分そのものである．

重なりで $s_\alpha=g_{\alpha\beta}s_\beta$ だから，$Ds$ が切断として整合する（$(Ds)_\alpha = g_{\alpha\beta}(Ds)_\beta$）条件は
\begin{equation}
  A_\beta = g_{\alpha\beta}^{-1}A_\alpha g_{\alpha\beta} + g_{\alpha\beta}^{-1}dg_{\alpha\beta}
  \label{eq:ch14-A-transform}
\end{equation}
である．実際 $(d+A_\alpha)(g_{\alpha\beta}s_\beta) = g_{\alpha\beta}\bigl(d + g_{\alpha\beta}^{-1}dg_{\alpha\beta} + g_{\alpha\beta}^{-1}A_\alpha g_{\alpha\beta}\bigr)s_\beta$．
同じ式は（同一の開集合上での）ゲージ変換 $s\mapsto g^{-1}s$ に対する
\begin{equation}
  A\mapsto A^g = g^{-1}Ag + g^{-1}dg
\end{equation}
とも読める．$U(1)$ で $g=e^{i\theta}$ とすれば $A^g = A + i\,d\theta$，すなわち $A^{\rm phys}\mapsto A^{\rm phys}-\frac1g d\theta$ である
（ゲージ変換の向きの違いを除いて電磁気学のゲージ変換）．

$A$ は非斉次項 $g^{-1}dg$ のために $\End E$ 値の1形式（テンソル）ではない．しかし二つの接続の差 $D-D'$ は
Leibniz 則の項が打ち消すので $\Omega^1(X,\End E)$ の元である．したがって接続全体 $\mathcal{A}$ は
$\Omega^1(X,\End E)$ を平行移動の群とする\textbf{アフィン空間}である（無限次元）．
接続は1の分割を用いて常に存在する．

共変微分は $E$ 値 $p$ 形式に $D(\omega\otimes s) = d\omega\otimes s + (-1)^p\omega\wedge Ds$ で延長され，
$D:\Omega^p(X,E)\to\Omega^{p+1}(X,E)$ を得る．また $E$ の接続は $E^*,\ E\otimes F,\ \End E$ などに自然に誘導される．
$\End E$ 上では $D\xi = d\xi + [A,\xi]$ である．

\subsection{曲率}

$d^2=0$ と違い，$D^2$ は一般に零ではない．しかし $D^2(fs) = f D^2 s$ となり（確かめよ），$D^2$ は
微分を含まない $C^\infty(X)$-線形な作用素，すなわち $\End E$ 値2形式の掛け算である．

\begin{mydef}{曲率}{ch14-curvature}
接続 $D$ の\textbf{曲率}とは $F=D^2\in\Omega^2(X,\End E)$ をいう．局所的には
\begin{equation}
  F = dA + A\wedge A,\qquad F_{\mu\nu} = \partial_\mu A_\nu - \partial_\nu A_\mu + [A_\mu,A_\nu]
\end{equation}
である．$F=0$ のとき接続は\textbf{平坦}であるという．
\end{mydef}

実際 $D^2s = (d+A)(ds + As) = dA\,s - A\wedge ds + A\wedge ds + A\wedge A\,s = (dA + A\wedge A)s$．
ゲージ変換に対しては \eqref{eq:ch14-A-transform} から
\begin{equation}
  F\mapsto g^{-1}Fg
\end{equation}
と随伴的（共変的）に変換する．\eqref{eq:ch14-convention} によって $F_{\mu\nu}=-igF^{\rm phys}_{\mu\nu}$ であり，
非可換ゲージ理論の場の強さ $F^{\rm phys}_{\mu\nu}=\partial_\mu A^{\rm phys}_\nu-\partial_\nu A^{\rm phys}_\mu-ig[A^{\rm phys}_\mu,A^{\rm phys}_\nu]$ を再現する．

\begin{myprop}{Bianchi 恒等式}{ch14-bianchi}
曲率は $DF := dF + [A,F] = 0$ をみたす．
\end{myprop}
\begin{proof}
$F=dA+A\wedge A$ より $dF = dA\wedge A - A\wedge dA = (F-A\wedge A)\wedge A - A\wedge(F-A\wedge A) = F\wedge A - A\wedge F$．
よって $dF + A\wedge F - F\wedge A = 0$．抽象的には $D(D^2)=D^2(D)$ という自明な等式の言い換えである．
\end{proof}
電磁気（$U(1)$）では $dF=0$ となり，Maxwell 方程式の半分 $\nabla\cdot\mathbf{B}=0$，$\nabla\times\mathbf{E}+\partial_t\mathbf{B}=0$ を与える．

\subsection{Yang--Mills理論の幾何学的記述}

$X$ を向き付けられた Riemann（あるいは Lorentz）多様体とし，Hodge スター $*$ を用いる．構造群 $G=U(n)$ または $SU(n)$ のとき，
$\mathfrak{g}$ 上の正定値内積 $\langle\xi,\eta\rangle=-\Tr(\xi\eta)$ を用いて \textbf{Yang--Mills 作用}を
\begin{equation}
  S_{\rm YM}[A] = -\frac{1}{2g^2}\int_X\Tr(F\wedge *F) = \frac{1}{4g^2}\int_X d^nx\sqrt{|\eta|}\,(-\Tr F_{\mu\nu}F^{\mu\nu})
\end{equation}
と定める．$F=-igF^{\rm phys}$ を代入すれば，$-\Tr F_{\mu\nu}F^{\mu\nu}=g^2\Tr F^{\rm phys}_{\mu\nu}F^{{\rm phys}\,\mu\nu}$ であり，
通常の $\frac14\int\Tr F^{\rm phys}_{\mu\nu}F^{{\rm phys}\,\mu\nu}$（規格化 $\Tr T^aT^b=\delta^{ab}$）に一致する．
$F\mapsto g^{-1}Fg$ とトレースの巡回性により，$S_{\rm YM}$ はゲージ不変である．

変分 $A\to A+\delta A$ で $\delta F = D\delta A$ だから，$\delta S \propto \int\Tr(D\delta A\wedge *F)$ を部分積分して
\textbf{Yang--Mills 方程式}
\begin{equation}
  D{*F} = 0
\end{equation}
を得る．Bianchi 恒等式 $DF=0$ と合わせて，非可換版の Maxwell 方程式の全体となる．

4次元 Euclid 空間では $*^2=1$ なので $F=F^++F^-$（$*F^\pm=\pm F^\pm$）と分解でき，
\begin{equation}
  S_{\rm YM} = -\frac{1}{2g^2}\int\Tr(F^+\wedge F^+ - F^-\wedge F^-)\ \geq\ \frac{1}{2g^2}\left|\int\Tr(F\wedge F)\right|
\end{equation}
が成り立つ（$F^\pm\wedge *F^\pm = \pm F^\pm\wedge F^\pm$，交差項 $\Tr(F^+\wedge F^-)$ は消え，$\int\Tr(F\wedge F)=\int\Tr(F^+\wedge F^+ + F^-\wedge F^-)$ であり，$-\Tr(F^+\wedge F^+)\geq0$，$\Tr(F^-\wedge F^-)\geq0$ による）．右辺は位相不変量（第2 Chern 数に比例，第15章）であり，等号は\textbf{（反）自己双対方程式} $F=\pm*F$ のとき成り立つ．
その解が\textbf{インスタントン}であり，Bianchi 恒等式から自動的に Yang--Mills 方程式を解く．
このように「作用の下限を与える1階の方程式」は BPS 方程式の典型であり，14.6.5節で再び現れる．

\subsection{曲率と局所的な幾何}

曲率の幾何学的意味は\textbf{ホロノミー}にある．点 $x$ から座標方向 $\mu,\nu$ に $\epsilon$ ずつ進む小さな正方形のループ $\gamma$ に沿った平行移動は
\begin{equation}
  P_\gamma = 1 - \epsilon^2 F_{\mu\nu}(x) + O(\epsilon^3)
\end{equation}
となる（符号はループの向きに依存）．すなわち曲率は「無限小ループに沿ったファイバーの回転」である．
有限のループに対しては経路順序積
$P_\gamma = \mathcal{P}\exp\left(-\oint_\gamma A\right)$
（Wilson ループ）が対応し，$U(1)$ ではStokes の定理から $P_\gamma = \exp(-\int_\Sigma F)$（$\partial\Sigma=\gamma$）である．
物理的には Aharonov--Bohm 位相 $\exp(ig\oint A^{\rm phys})=\exp(ig\Phi_B)$ である．

\begin{myprop}{平坦接続とホロノミー表現}{ch14-flat-holonomy}
$X$ を連結多様体，$D$ を平坦接続とすると，ホロノミーはループのホモトピー類のみに依存し，表現
\begin{equation}
  \mathrm{hol}_D:\pi_1(X,x_0)\to G
\end{equation}
を与える．逆に表現 $\rho:\pi_1(X)\to G$ から，普遍被覆 $\tilde X$ を用いて $\tilde X\times_\rho G$ という平坦接続つきの束が得られる．
\end{myprop}
\begin{proof}[証明の概略]
平坦なら局所的に $A=h^{-1}dh$ と書け（Frobenius の定理），ゲージ変換で $A=0$ にできる．したがって局所的にはホロノミーは自明であり，
ループの小さな変形でホロノミーは変わらない．よってホモトピー不変である．
\end{proof}
単連結な空間上の平坦接続は自明だが，$\pi_1$ が非自明な空間（トーラス，Riemann 面，Aharonov--Bohm のソレノイドの外側）では，
平坦接続でも観測可能なホロノミー（Wilson ライン）をもつ．これが 14.6.3節の平坦接続のモジュライの出発点である．

\begin{physbox}[曲率＝場の強さ，ホロノミー＝Wilson ループ]
\begin{center}
\begin{tabular}{c|c}
物理 & 幾何 \\ \hline
ゲージポテンシャル $A_\mu$ & 接続形式 \\
場の強さ $F_{\mu\nu}$ & 曲率 \\
ゲージ変換 & 束の自己同型 \\
物質場 & 同伴ベクトル束の切断 \\
Wilson ループ & ホロノミー \\
磁荷・インスタントン数 & 特性類（第15章）
\end{tabular}
\end{center}
\end{physbox}

\exprob{ch14-pr-curv}
ゲージ変換 $A\mapsto g^{-1}Ag+g^{-1}dg$ のもとで $F=dA+A\wedge A$ が $g^{-1}Fg$ に変換することを，直接計算で示せ．

\exprob{ch14-pr-mono}
$S^2$ を北側 $U_N=\{\theta\neq\pi\}$ と南側 $U_S=\{\theta\neq0\}$ で覆い，局所接続（ユニタリ枠での $\mathfrak{u}(1)$ 値1形式）
$A_N = -\frac{in}{2}(1-\cos\theta)d\varphi$，$A_S=\frac{in}{2}(1+\cos\theta)d\varphi$ をとる．
$A_N,A_S$ がそれぞれ $U_N,U_S$ 上で滑らかであること，および例\ref{ex:ch14-dirac} の遷移関数 $g_{NS}=e^{in\varphi}$
（$s_S=s_Ng_{NS}$，すなわち \eqref{eq:ch14-A-transform} で $\alpha=N,\beta=S$）について \eqref{eq:ch14-A-transform} の関係にあることを確かめ，
$S^2$ に $d\theta\wedge d\varphi$ が正となる向きを入れて $\frac{i}{2\pi}\int_{S^2}F=n$ を示せ．
（この規約で $\frac{i}{2\pi}\int F$ は磁荷 $n$ に等しく，例\ref{ex:ch14-o-1} の $\OO(-1)$ は $n=-1$ にあたる．）

% ----------------------------------------------------------
\section{複素ベクトル束}
% ----------------------------------------------------------

\subsection{複素ベクトル束}

ここからは底空間 $X$ を $n$ 次元複素多様体とする．ファイバーが複素ベクトル空間で，遷移関数が $\GL(r,\C)$ 値の滑らかな関数である束を
\textbf{複素ベクトル束}という（底空間の複素構造とは無関係な概念である）．
$E$ 値の微分形式は $(p,q)$ 型に分解し，$\Omega^{p,q}(X,E)$ と書く．接続も
\begin{equation}
  D = D^{1,0} + D^{0,1},\qquad D^{0,1}:\Omega^{p,q}(X,E)\to\Omega^{p,q+1}(X,E)
\end{equation}
と分解する．局所的には $A = A^{1,0}+A^{0,1}$，$A^{1,0}=A_i dz^i$，$A^{0,1}=A_{\bar i}d\bar z^{\bar i}$ である．
同様に曲率も
\begin{equation}
  F = F^{2,0} + F^{1,1} + F^{0,2}
\end{equation}
と分解する．$D^{0,1}$ は $\bar\partial$ の「ねじれた」版であり，以下で見るように，これが正則構造を担う．

\subsection{正則ベクトル束}

\begin{mydef}{正則ベクトル束}{ch14-holvb}
複素多様体 $X$ 上の複素ベクトル束 $E$ が\textbf{正則ベクトル束}であるとは，遷移関数 $g_{\alpha\beta}:U_\alpha\cap U_\beta\to\GL(r,\C)$ が
すべて正則であるような局所自明化が与えられていることをいう．このとき切断 $s$ が\textbf{正則}であるとは，各局所成分 $s_\alpha$ が正則であることをいう．
正則切断全体を $H^0(X,E)$ と書く．
\end{mydef}

遷移関数が正則なので，$s_\alpha=g_{\alpha\beta}s_\beta$ の正則性は自明化によらない．同じ理由で，作用素
\begin{equation}
  \bar\partial_E:\Omega^{p,q}(X,E)\to\Omega^{p,q+1}(X,E),\qquad (\bar\partial_E s)_\alpha := \bar\partial s_\alpha
\end{equation}
が well-defined となる（$\bar\partial(g_{\alpha\beta}s_\beta) = g_{\alpha\beta}\bar\partial s_\beta$ だから）．
$\bar\partial_E$ は Leibniz 則 $\bar\partial_E(fs)=\bar\partial f\otimes s + f\bar\partial_E s$ と $\bar\partial_E^2=0$ をみたし，
正則切断は $\bar\partial_E s=0$ の解である．Dolbeault コホモロジー $H^{p,q}_{\bar\partial}(X,E)$ は Dolbeault の定理により
層コホモロジー $H^q(X,\Omega^p_X\otimes\OO(E))$ に一致する（第11章・第15章）．

\begin{myex}{正則ベクトル束の例}{ch14-holvb-ex}
(1) 正則接束 $T^{1,0}X$：遷移関数は座標変換の Jacobi 行列 $(\partial w^i/\partial z^j)$ で正則．
(2) 正則余接束 $\Omega^1_X=(T^{1,0}X)^*$ とその外積 $\Omega^p_X$，標準束 $K_X=\Omega^n_X$．
(3) $\PP^n$ 上の $\OO(d)$（例\ref{ex:ch14-taut}）．
(4) 複素トーラス $\C/\Lambda$ 上で，因子 $\lambda$ をもつ正則関数 $e_\lambda(z)$ による同一視 $(z,v)\sim(z+\lambda,e_\lambda(z)v)$ から得られる直線束．$\Lambda=\Z+\Z\tau$，$e_1=1$，$e_\tau = e^{-2\pi i z - \pi i\tau}$ とすると，正則切断は Jacobi テータ関数 $\vartheta(z,\tau)=\sum_{n\in\Z}e^{\pi i n^2\tau+2\pi i n z}$ である．
\end{myex}

同じ複素ベクトル束（$C^\infty$ 束として）の上に，互いに同型でない正則構造が多数ありうる．
たとえば種数 $g$ の Riemann 面上の次数0の正則直線束全体は $g$ 次元の複素トーラス（Jacobi 多様体）$\mathrm{Jac}(\Sigma)=\Pic^0(\Sigma)$
をなすが，$C^\infty$ 束としてはすべて自明である．正則構造の分類がモジュライ問題（14.6節）である．

\subsection{Hermitian構造}

\begin{mydef}{Hermitian 計量}{ch14-hermitian}
複素ベクトル束 $E$ 上の \textbf{Hermitian 計量} $h$ とは，各ファイバー $E_x$ 上の Hermite 内積 $h_x$ で，$x$ について滑らかなものをいう．
局所枠 $e_i$ について $H_{ji}=h(e_i,e_j)$ とおくと，$H$ は正定値 Hermite 行列値関数で，$h(s,t)=t^\dagger H s$ である．
接続 $D$ が $h$ と\textbf{両立する}（ユニタリ接続）とは
\begin{equation}
  d\,h(s,t) = h(Ds,t) + h(s,Dt)\qquad(s,t\in\Gamma(X,E))
\end{equation}
をみたすことをいう．
\end{mydef}

1の分割により Hermitian 計量は常に存在する．$h$ を決めることは，構造群を $\GL(r,\C)$ から $U(r)$ へ簡約することと同値である
（$h$ に関する正規直交枠を局所枠にとれば遷移関数はユニタリ）．ユニタリ枠で見ると，両立条件は $A^\dagger=-A$，すなわち $A$ が $\mathfrak{u}(r)$ 値であることを意味し，
曲率も $\mathfrak{u}(r)$ 値（$F^\dagger=-F$）となる．物理のゲージ理論で $U(r)$ をゲージ群とするとき，暗黙のうちに Hermitian 計量が固定されている．

一方，正則構造は複素化された群 $\GL(r,\C)$ の言葉であり，正則枠と正規直交枠は一般に一致しない．
この「二つの構造（正則構造と Hermitian 構造）」をつなぐのが次の Chern 接続である．

\subsection{正則構造と接続}

\begin{mythm}{Chern 接続}{ch14-chern-conn}
$(E,h)$ を Hermitian 計量をもつ正則ベクトル束とする．このとき
\begin{equation}
  \text{(i) } D \text{ は } h \text{ と両立する},\qquad \text{(ii) } D^{0,1} = \bar\partial_E
\end{equation}
をみたす接続 $D$ が一意的に存在する．これを \textbf{Chern 接続}という．正則枠に関して
\begin{equation}
  A = H^{-1}\partial H,\qquad F = \bar\partial(H^{-1}\partial H)
  \label{eq:ch14-chern}
\end{equation}
であり，曲率は $(1,1)$ 型である：$F^{2,0}=F^{0,2}=0$．
\end{mythm}

\begin{proof}
正則枠 $e_i$ をとる．(ii) は $D^{0,1}e_i = \bar\partial_E e_i = 0$，すなわち $A$ が $(1,0)$ 型であることと同値である．
(i) を $s=e_i,t=e_j$ に適用すると，行列で $dH = A^\dagger H + HA$．$A$ が $(1,0)$ 型なら $A^\dagger$ は $(0,1)$ 型なので，
$(1,0)$ 部分をとって $\partial H = HA$，すなわち $A=H^{-1}\partial H$ と一意に決まる．
このとき $(0,1)$ 部分 $\bar\partial H = A^\dagger H = (\bar\partial H)H^{-1}H$ は自動的に成り立つので存在も示された．
（一意性から，枠の取り替えに対して整合的に貼り合わさる．）
曲率は $F = dA + A\wedge A = \bar\partial A + (\partial A + A\wedge A)$ であり，
$\partial(H^{-1}\partial H) = -H^{-1}\partial H\wedge H^{-1}\partial H = -A\wedge A$ なので $F=\bar\partial A$ となる．
\end{proof}

\begin{myex}{$\OO(-1)$ の Chern 接続}{ch14-o-1}
$\PP^1$ 上の $\OO(-1)$ に $\C^2$ の標準内積から計量を誘導する．$U_0$ 上の正則枠 $e_0=(1,z)$ について $H=|e_0|^2=1+|z|^2$．
よって $A = \partial\log(1+|z|^2)=\frac{\bar z\,dz}{1+|z|^2}$，
\begin{equation}
  F = \bar\partial\partial\log(1+|z|^2) = -\frac{dz\wedge d\bar z}{(1+|z|^2)^2}.
\end{equation}
$dz\wedge d\bar z=-2i\,dx\wedge dy$ を用いると $\frac{i}{2\pi}\int_{\PP^1}F = -\frac{1}{\pi}\int_{\R^2}\frac{dx\,dy}{(1+r^2)^2}=-1$．
この整数が $\OO(-1)$ の次数（第1 Chern 数）であり，Dirac モノポールの磁荷 $-1$ に対応する．Chern 類の一般論は第15章で扱う．
\end{myex}

定理\ref{thm:ch14-chern-conn} の逆が成り立つことが，複素幾何とゲージ理論をつなぐ鍵である．

\begin{mythm}{接続による正則構造の特徴づけ}{ch14-koszul-malgrange}
$E$ を複素多様体 $X$ 上の $C^\infty$ 複素ベクトル束とし，$D$ をその接続とする．
$F^{0,2}=0$（すなわち $(D^{0,1})^2=0$）ならば，$E$ 上の正則構造で $\bar\partial_E = D^{0,1}$ となるものが一意的に存在する．
\end{mythm}
\begin{proof}[証明の概略]
局所的に $D^{0,1} = \bar\partial + A^{0,1}$ であり，$(D^{0,1})^2 = \bar\partial A^{0,1}+A^{0,1}\wedge A^{0,1}=F^{0,2}$．
$F^{0,2}=0$ のとき，各点の近傍で $\bar\partial g = -A^{0,1}g$ をみたす $\GL(r,\C)$ 値関数 $g$ が存在する（Newlander--Nirenberg 型の可積分性定理，
$\dim_\C X=1$ なら $F^{0,2}=0$ は自動的で，$\bar\partial$ 方程式を解けばよい）．
このとき $e_i' = e_j g_{ji}$ は $D^{0,1}e'_i=0$ をみたす局所枠であり，二つのそのような枠の間の遷移関数は $\bar\partial$ で消える，すなわち正則である．
\end{proof}

すなわち，$C^\infty$ 束 $E$ を固定したとき
\begin{align}
  \{E\text{ 上の正則構造}\} &\;\longleftrightarrow\; \{\bar\partial\text{ 作用素 } \bar\partial_E \text{ で } \bar\partial_E^2=0\} \notag\\
  &\;\longleftrightarrow\; \{\text{ユニタリ接続 } A \text{ で } F^{0,2}=0\}
\end{align}
（最後の対応は Hermitian 計量 $h$ を固定したとき，Chern 接続を通じて）．複素化ゲージ群 $\mathcal{G}_\C$（$\GL(r,\C)$ 値ゲージ変換）は
$\bar\partial_E\mapsto g^{-1}\bar\partial_E g$ によって正則構造に作用し，その軌道が正則束の同型類である．

\subsection{複素幾何とゲージ理論}

ユニタリ接続のうち $F^{0,2}=0$ をみたすものは正則構造と同じであった．では正則構造（$\mathcal{G}_\C$ 軌道）の中で，
ゲージ理論的に「最良の」接続を選ぶ条件は何か．$X$ を Kähler 形式 $\omega$ をもつコンパクト Kähler 多様体とする（第15章）．
$\Lambda_\omega$ を $\omega$ との縮約（$(1,1)$ 形式から関数への写像）とする．

\begin{mydef}{Hermitian--Yang--Mills 方程式}{ch14-hym}
Hermitian 計量つき正則ベクトル束 $(E,h)$ の Chern 接続が
\begin{equation}
  F^{0,2}=0\ (\text{自動的}),\qquad i\Lambda_\omega F = \lambda\cdot\mathbf{1}_E
  \label{eq:ch14-hym}
\end{equation}
（$\lambda$ は定数）をみたすとき，$h$ を \textbf{Hermitian--Einstein 計量}，接続を \textbf{Hermitian--Yang--Mills（HYM）接続}という．
局所座標では $g^{i\bar j}F_{i\bar j}\propto\mathbf{1}$ である．
\end{mydef}

定数 $\lambda$ は位相的に決まる．$\deg E := \int_X c_1(E)\wedge\omega^{n-1}$（$c_1(E)=[\frac{i}{2\pi}\Tr F]$，第15章）とし，
\textbf{傾き}を $\mu(E)=\deg E/\rank E$ とすると，$\lambda$ は $\mu(E)$ に比例する．
HYM 接続は Yang--Mills 方程式 $D{*F}=0$ をみたす（Kähler 恒等式と Bianchi 恒等式による）．
$\dim_\C X=2$ では HYM 方程式は（$\lambda=0$ のとき）反自己双対方程式 $F^+=0$ と同値である（$F^{0,2}=0$ と $\Lambda F=0$ で $F^+$ の3成分）．

\begin{mydef}{傾き安定性}{ch14-stability}
コンパクト Kähler 多様体上の正則ベクトル束 $E$ が\textbf{安定}（stable）であるとは，$0<\rank\sheaf{F}<\rank E$ をみたす任意の連接部分層 $\sheaf{F}\subset E$ に対し
\begin{equation}
  \mu(\sheaf{F}) < \mu(E)
\end{equation}
が成り立つことをいう．$\leq$ のとき半安定，安定束の直和で同じ傾きのものを多重安定（polystable）という．
\end{mydef}

\begin{mythm}{Donaldson--Uhlenbeck--Yau（小林--Hitchin 対応）}{ch14-duy}
コンパクト Kähler 多様体上の既約な正則ベクトル束 $E$ が Hermitian--Einstein 計量をもつための必要十分条件は，$E$ が安定であることである．
このとき Hermitian--Einstein 計量は定数倍を除いて一意である．（一般に，多重安定 $\Leftrightarrow$ HYM 接続をもつ．）
\end{mythm}

「HYM $\Rightarrow$ 多重安定」は比較的やさしく（小林，Lübke），部分束の第二基本形式を用いた曲率の比較で示せる．
「安定 $\Rightarrow$ HYM」が難しい方向であり，Riemann 面では Narasimhan--Seshadri（Donaldson による別証明），
代数曲面では Donaldson，一般の Kähler 多様体では Uhlenbeck--Yau による．
この定理は\textbf{解析的な条件（偏微分方程式の可解性）と代数的な条件（部分層の次数の不等式）の同値性}を述べており，
GAGA（第13章）と並ぶ「解析と代数の橋」である．

\begin{myex}{$\PP^1$ 上の安定性}{ch14-p1-stab}
Grothendieck の定理により $\PP^1$ 上の正則ベクトル束は $\OO(a_1)\oplus\cdots\oplus\OO(a_r)$ と分解する．
$\OO(a)\oplus\OO(b)$（$a\geq b$）は部分束 $\OO(a)$ の傾き $a$ と全体の傾き $(a+b)/2$ を比べて，$a=b$ のとき多重安定（半安定）で，$a>b$ なら不安定である．
実際，$\OO(a)\oplus\OO(a)$ には Fubini--Study 計量から誘導される HYM 接続（各成分が定曲率）が存在する．
\end{myex}

\begin{physbox}[HYM 方程式と超対称性]
ヘテロ型弦理論を Calabi--Yau 多様体 $X$ 上でコンパクト化するとき，内部空間のゲージ場は4次元で $\mathcal{N}=1$ 超対称性を保つために
$F^{2,0}=F^{0,2}=0$，$\Lambda_\omega F = 0$ をみたさねばならない（gaugino の超対称変換 $\delta\chi = F_{mn}\Gamma^{mn}\epsilon=0$ から）．
すなわち HYM 接続であり，DUY 定理により「安定な正則ベクトル束を代数幾何的に構成する」問題に帰着する．
これが弦理論による標準模型の構成で代数幾何（モナド構成，スペクトル被覆）が使われる理由である．
4次元の言葉では $F^{0,2}=0$ は F 項条件，$\Lambda F=0$ は D 項条件にあたる（14.6.5節）．
\end{physbox}

\exprob{ch14-pr-chern}
直線束 $L$ の Hermitian 計量を正則枠で $h=|e|^2=e^{-\varphi}$ と書くとき，Chern 接続と曲率が $A=-\partial\varphi$，$F=\partial\bar\partial\varphi$ となることを示せ．

\exprob{ch14-pr-stab}
Riemann 面上の正則直線束は常に安定であることを示せ．また $\PP^1$ 上の $\OO(1)\oplus\OO(-1)$ が不安定であることを示せ．

% ----------------------------------------------------------
\section{Higgs束}
% ----------------------------------------------------------

\subsection{Higgs場の幾何学的背景}

ゲージ理論の Higgs 場は随伴表現（あるいは他の表現）に値をとるスカラー場であり，幾何学的には同伴束の切断である．
随伴表現の Higgs 場は $\End E$（$\mathfrak{u}(r)$ に値をとる部分）の切断で，共変微分 $D\Phi = d\Phi+[A,\Phi]$ をもつ．

Higgs 束の動機は\textbf{次元還元}にある．4次元 Euclid 空間 $\R^4$（座標 $x^1,\dots,x^4$）上の反自己双対方程式 $F^+=0$ において，
場が $x^3,x^4$ によらないと仮定すると，$A_3,A_4$ は2次元 $\R^2=\C$（$z=x^1+ix^2$）上の随伴スカラー場となる．
\begin{equation}
  \Phi := \tfrac12(A_3 - iA_4)\,dz
\end{equation}
とまとめると，$F^+=0$ の3成分は
\begin{equation}
  F_A + [\Phi,\Phi^*] = 0,\qquad \bar\partial_A\Phi = 0
  \label{eq:ch14-hitchin-pre}
\end{equation}
という2次元の方程式（Hitchin 方程式）に化ける（$\Phi^*$ は $h$ に関する随伴をとり形式の部分は複素共役したもの）．
注目すべきは，$\Phi$ が\textbf{1形式}（$dz$ に比例）として現れ，正則性 $\bar\partial_A\Phi=0$ をみたすことである．
Riemann 面に一般化すれば，$\Phi$ は $\End E\otimes K_\Sigma$ の正則切断となる．

\subsection{Higgs束の定義}

\begin{mydef}{Higgs 束}{ch14-higgs}
複素多様体 $X$ 上の \textbf{Higgs 束}とは，正則ベクトル束 $E$ と正則切断
\begin{equation}
  \Phi\in H^0(X,\End E\otimes\Omega^1_X)
\end{equation}
で $\Phi\wedge\Phi=0$ をみたすものの組 $(E,\Phi)$ をいう．$\Phi$ を \textbf{Higgs 場}という．
$X=\Sigma$ が Riemann 面のときは $\Omega^1_\Sigma=K_\Sigma$ であり，$\Phi\wedge\Phi=0$ は自動的にみたされるので，
\begin{equation}
  \Phi\in H^0(\Sigma,\End E\otimes K_\Sigma)
\end{equation}
を与えればよい．局所的には $\Phi = \phi(z)\,dz$，$\phi$ は正則な行列値関数である．
\end{mydef}

二つの Higgs 束 $(E,\Phi),(E',\Phi')$ の同型とは，正則同型 $f:E\to E'$ で $\Phi' f = f\Phi$ をみたすものをいう．
$\Phi$ は正則な切断なので，複素ゲージ変換に対して $\Phi\mapsto g^{-1}\Phi g$ と変換する．

\begin{mydef}{Higgs 束の安定性}{ch14-higgs-stab}
Riemann 面上の Higgs 束 $(E,\Phi)$ が\textbf{安定}であるとは，$\Phi$ で不変な（$\Phi(F)\subset F\otimes K$）任意の真の部分束 $0\neq F\subsetneq E$ に対し
$\mu(F)<\mu(E)$ が成り立つことをいう．
\end{mydef}
$\Phi$ 不変な部分束だけを検査するので，$E$ 自身が不安定でも $(E,\Phi)$ は安定になりうる．

\subsection{Higgs場と正則構造}

Higgs 束に Hermitian 計量 $h$ を入れ，$E$ の Chern 接続 $A$（14.4.4節）を考える．次の方程式系を\textbf{Hitchin 方程式}（自己双対方程式）という：
\begin{equation}
  F_A + [\Phi,\Phi^*] = 0,\qquad \bar\partial_A\Phi = 0.
  \label{eq:ch14-hitchin}
\end{equation}
第2式は「$\Phi$ が正則」というだけであり，Higgs 束の定義に含まれている．第1式が計量 $h$ に対する方程式であり，
HYM 方程式 $i\Lambda F=\lambda$ の Higgs 版である．（$\deg E\neq0$ のときは右辺に定数 $\propto\mu(E)\,\omega\,\mathbf{1}$ を加える．）

\begin{mythm}{Hitchin，Simpson}{ch14-hitchin-thm}
コンパクト Riemann 面（種数 $g\geq2$）上の Higgs 束 $(E,\Phi)$ が Hitchin 方程式の解となる計量をもつための必要十分条件は，
$(E,\Phi)$ が多重安定であることである．
\end{mythm}

これは DUY 定理の Higgs 版である．さらに $\nabla := D_A + \Phi + \Phi^*$ とおくと，Hitchin 方程式は $\nabla$ が\textbf{平坦な $\GL(r,\C)$ 接続}であることと
（調和性条件の下で）同値になる．実際，曲率を $(1,1)$ 成分で計算すると
$F_\nabla = F_A + [\Phi,\Phi^*] + \bar\partial_A\Phi + \partial_A\Phi^*$ となり，\eqref{eq:ch14-hitchin} のもとで $0$ である．
こうして
\begin{center}
安定 Higgs 束 $\longleftrightarrow$ Hitchin 方程式の解 $\longleftrightarrow$ 平坦 $\GL(r,\C)$ 接続（既約）
\end{center}
という対応（\textbf{非可換 Hodge 対応}，Hitchin--Donaldson--Corlette--Simpson）が得られる．
$U(1)$ の場合は，平坦 $\C^*$ 接続（$H^1(\Sigma,\C)$）が正則直線束（$H^{0,1}$）と正則1形式（$H^{1,0}$）に分かれる Hodge 分解
$H^1(\Sigma,\C)=H^{1,0}\oplus H^{0,1}$ そのものであり，非可換 Hodge 対応はその非可換版である．

\subsection{Higgs束の基本的な例}

\begin{myex}{階数1}{ch14-higgs-rank1}
$E=L$ が直線束なら $\End L=\OO$ なので，Higgs 場は正則1形式 $\Phi\in H^0(\Sigma,K_\Sigma)\cong\C^g$ である．
Hitchin 方程式は $F_A=0$（$[\Phi,\Phi^*]=0$）となり，$L$ は平坦 $U(1)$ 束．
次数0の Higgs 直線束全体は $\mathrm{Jac}(\Sigma)\times H^0(K)\cong T^*\mathrm{Jac}(\Sigma)$ である
（$\mathrm{Jac}$ の余接空間は $H^0(K)$）．
\end{myex}

\begin{myex}{Hitchin 切断（一意化 Higgs 束）}{ch14-higgs-uniform}
$\Sigma$ の種数を $g\geq2$ とし，$K^{1/2}$ を $K$ の平方根（スピン構造）とする．
\begin{equation}
  E = K^{1/2}\oplus K^{-1/2},\qquad \Phi = \begin{pmatrix} 0 & q\\ 1 & 0\end{pmatrix},\quad q\in H^0(\Sigma,K^2)
\end{equation}
とおく．ここで左下の $1$ は $\Hom(K^{1/2},K^{-1/2})\otimes K = \OO$ の定数切断，右上の $q$ は $\Hom(K^{-1/2},K^{1/2})\otimes K=K^2$ の切断（2次微分）である．
$E$ 自身は不安定（部分束 $K^{1/2}$ の次数 $g-1>0=\mu(E)$）だが，$K^{1/2}$ は $\Phi$ 不変でなく（左下成分 $1$ が $K^{1/2}$ を $K^{-1/2}\otimes K$ へ写す），
$\Phi$ 不変な直線部分束はすべて次数が負であることが示せるため，
$(E,\Phi)$ は安定である．$q=0$ のとき Hitchin 方程式の解は $\Sigma$ の定曲率 $-1$ の双曲計量（一意化定理，第13章）を与え，
対応する平坦 $\SL(2,\R)$ 接続は $\pi_1(\Sigma)\subset\mathrm{PSL}(2,\R)$ の Fuchs 表現である．
\end{myex}

\begin{myex}{$\PP^1$ と楕円曲線}{ch14-higgs-p1}
$\PP^1$ では $K=\OO(-2)$ であり，半安定な $E=\OO(a)^{\oplus r}$ に対して $H^0(\End E\otimes K)=H^0(\OO(-2))^{\oplus r^2}=0$．
すなわち非自明な Higgs 場はない．楕円曲線 $E_\tau$ では $K\cong\OO$ なので Higgs 場は $\End E$ の正則切断そのものである．
種数 $g\geq 2$ ではじめて Higgs 束の理論が豊かになる．
\end{myex}

\subsection{Higgs束とモジュライ}

Higgs 束の同型類の空間には，Higgs 場の不変量を取り出す写像がある．$\Phi$ の特性多項式
\begin{equation}
  \det(x - \Phi) = x^r + a_1 x^{r-1} + \cdots + a_r,\qquad a_k = (-1)^k\Tr\textstyle\bigwedge^k\Phi\in H^0(\Sigma,K^k)
\end{equation}
の係数は $\Phi\mapsto g^{-1}\Phi g$ で不変である（$\Phi$ は $K$ に値をとるので $a_k$ は $K^k$ の切断）．
これにより
\begin{equation}
  h:\mathcal{M}_{\rm Higgs}\to\mathcal{B}=\bigoplus_{k=1}^{r}H^0(\Sigma,K^k),\qquad (E,\Phi)\mapsto(a_1,\dots,a_r)
\end{equation}
が定まる．これを \textbf{Hitchin fibration} という．$\SL(r)$ の場合は $a_1=0$ とする．$\SL(2)$ なら $\mathcal{B}=H^0(K^2)$ で，
$\dim\mathcal{B} = 3g-3$（Riemann--Roch）．これはモジュライ空間 $\mathcal{M}_{\rm Higgs}$ の次元 $6g-6$ のちょうど半分である．

$a\in\mathcal{B}$ に対し，$K$ の全空間 $\mathrm{Tot}(K)$ の中の曲線
\begin{equation}
  \Sigma_a = \{x\in\mathrm{Tot}(K):x^r + a_1x^{r-1}+\cdots+a_r=0\}
\end{equation}
を\textbf{スペクトル曲線}という．これは $\Sigma$ の $r$ 重分岐被覆であり，その点は $\Phi$ の固有値（固有1形式）である．
一般の $a$ に対し，ファイバー $h^{-1}(a)$ はスペクトル曲線上の直線束のモジュライ，すなわち $\Sigma_a$ の Jacobi 多様体（の一部）であり，
コンパクトな複素トーラス（アーベル多様体）となる．$\mathcal{M}_{\rm Higgs}$ は正則シンプレクティック多様体で，
$h$ の成分は互いに Poisson 可換な関数をなす．すなわち $h$ は\textbf{代数的完全可積分系}である（Hitchin 系）．

\begin{figure}[ht]
\centering
\begin{tikzpicture}[scale=1.0]
  \draw[thick] (0,0) -- (7,0) node[right]{$\mathcal{B}=\bigoplus H^0(K^k)$};
  \foreach \x in {1,2.5,4,5.5} {
    \draw[thick,blue!70!black] (\x,1.8) ellipse (0.45 and 0.9);
    \draw[blue!70!black] (\x-0.12,1.85) .. controls (\x-0.05,1.7) and (\x+0.05,1.7) .. (\x+0.12,1.85);
    \draw[blue!70!black] (\x-0.09,1.78) .. controls (\x-0.03,1.86) and (\x+0.03,1.86) .. (\x+0.09,1.78);
    \fill (\x,0) circle (1.3pt);
  }
  \draw[thick,red!70!black] (6.5,1.8) .. controls (6.1,2.7) and (6.9,2.7) .. (6.5,1.8) .. controls (6.1,0.9) and (6.9,0.9) .. (6.5,1.8);
  \fill (6.5,0) circle (1.3pt);
  \node[red!70!black,right] at (6.8,1.8) {\small 特異ファイバー};
  \draw[->,thick] (-0.5,1.8) -- (-0.5,0.3) node[midway,left]{$h$};
  \node at (-0.5,2.4) {$\mathcal{M}_{\rm Higgs}$};
  \node[blue!70!black] at (3.25,3.0) {\small 一般ファイバー：Jacobi 多様体（トーラス）};
\end{tikzpicture}
\caption{Hitchin fibration の模式図．一般のファイバーはスペクトル曲線の Jacobi 多様体であり，判別式の上で退化する．}
\label{fig:ch14-hitchin}
\end{figure}

\begin{physbox}[Hitchin 系と Seiberg--Witten 理論]
4次元 $\mathcal{N}=2$ 超対称ゲージ理論（特に Gaiotto の class $\mathcal{S}$ 理論）を円周上にコンパクト化すると，
3次元の低エネルギー理論の真空のモジュライ空間が Hitchin 系のモジュライ空間 $\mathcal{M}_{\rm Higgs}$ になる．
底空間 $\mathcal{B}$ は4次元のクーロン枝，スペクトル曲線は Seiberg--Witten 曲線，ファイバーのトーラスは電気・磁気の Wilson ラインの値である．
また Kapustin--Witten は，$\mathcal{N}=4$ 超対称 Yang--Mills 理論の S 双対性が，ゲージ群 $G$ と Langlands 双対群 ${}^LG$ の Hitchin 系の
ファイバーの双対性（SYZ 型のミラー対称性）として実現され，幾何学的 Langlands 対応につながることを示した．
\end{physbox}

\exprob{ch14-pr-hitchin-dim}
Riemann--Roch の定理 $\dim H^0(L)-\dim H^1(L)=\deg L+1-g$ と，$\deg L>2g-2$ のとき $H^1(L)=0$ であることを用いて，
$g\geq2$ のとき $\dim H^0(\Sigma,K^k) = (2k-1)(g-1)$（$k\geq2$）を示せ．これから $\SL(r)$ Hitchin 系の底空間の次元が $(r^2-1)(g-1)$ となることを導け．

% ----------------------------------------------------------
\section{モジュライ空間への接続}
% ----------------------------------------------------------

\subsection{ゲージ同値な構造の分類}

ゲージ理論の物理的配位は接続のゲージ同値類であり，さらに運動方程式や BPS 方程式の解に制限すると，
しばしば\textbf{有限次元}の空間が残る．一般に，
\begin{equation}
  \mathcal{M} = \frac{\{\text{方程式をみたす場の配位}\}}{\text{ゲージ群}\ \mathcal{G}}
\end{equation}
を\textbf{モジュライ空間}という．その次元は「方程式の線形化の解の数 $-$ ゲージ方向の数」として，楕円型複体の指数で計算される．
たとえば Riemann 面上の平坦接続では，変形複体
\begin{equation}
  \Omega^0(\mathfrak{g}_P)\xrightarrow{\ D\ }\Omega^1(\mathfrak{g}_P)\xrightarrow{\ D\ }\Omega^2(\mathfrak{g}_P)
\end{equation}
（第1の矢印は無限小ゲージ変換，第2は平坦性の線形化 $\delta F = D\delta A$）のコホモロジー $H^1$ が接空間となる．
モジュライ空間の一般論（変形理論，局所構造）は第16章で扱う．ここでは束・接続・Higgs 束に特有の事情をまとめる．

モジュライ空間を構成するときの典型的な問題は，ゲージ群の作用が自由でない点（自己同型をもつ点，可約接続）や，
同値関係が分離しない点（非 Hausdorff 性）が存在することである．安定性条件はこれを除くための条件であり，
安定な対象の自己同型はスカラー倍に限る（Schur の補題の類似）ため，モジュライ空間が良い空間になる．

\subsection{束のモジュライ}

種数 $g$ の Riemann 面 $\Sigma$ 上の階数 $r$，次数 $d$ の安定な正則ベクトル束の同型類の空間 $\mathcal{N}(r,d)$ は，
GIT（幾何学的不変式論）により準射影代数多様体として構成される（Mumford，Seshadri）．その次元は，
$E$ の無限小変形が $H^1(\Sigma,\End E)$ で，安定性から $H^0(\End E)=\C$（スカラー）であることを用いると，Riemann--Roch より
\begin{equation}
  \dim_\C\mathcal{N}(r,d) = \dim H^1(\End E) = r^2(g-1)+1
\end{equation}
である．$r=1$ では $\mathcal{N}(1,d)=\Pic^d(\Sigma)\cong\mathrm{Jac}(\Sigma)$（$g$ 次元複素トーラス）であり，
$\gcd(r,d)=1$ なら $\mathcal{N}(r,d)$ はコンパクトで滑らかな射影多様体になる．
行列式を固定した $\SL$ 版では次元は $(r^2-1)(g-1)$ である．

\subsection{接続のモジュライ}

Riemann 面 $\Sigma$ 上の平坦 $U(r)$ 接続のモジュライ空間を考える．命題\ref{prop:ch14-flat-holonomy} により平坦接続はホロノミー表現で決まり，
\begin{equation}
  \mathcal{M}_{\rm flat}(\Sigma,U(r)) = \Hom(\pi_1(\Sigma),U(r))/U(r)
\end{equation}
（$U(r)$ は共役で作用）となる．これを\textbf{表現の多様体}（指標多様体）という．
$\pi_1(\Sigma) = \langle a_1,b_1,\dots,a_g,b_g\mid\prod_{i=1}^g[a_i,b_i]=1\rangle$ であるから，具体的には
\begin{equation}
  \Hom(\pi_1(\Sigma),G) = \Bigl\{(A_1,B_1,\dots,A_g,B_g)\in G^{2g}:\prod_{i=1}^g A_iB_iA_i^{-1}B_i^{-1}=1\Bigr\}
\end{equation}
であり，$G=SU(r)$ なら実次元は $2g\dim G - \dim G - \dim G = (2g-2)\dim G$，複素次元にして $(r^2-1)(g-1)$ である．

\begin{mythm}{Narasimhan--Seshadri}{ch14-ns}
コンパクト Riemann 面 $\Sigma$（$g\geq 2$）上で，次数0の安定な正則ベクトル束の同型類と，
$\pi_1(\Sigma)$ の既約ユニタリ表現の共役類は一対一に対応する：
\begin{equation}
  \mathcal{N}^{s}(r,0) \;\cong\; \Hom^{\rm irr}(\pi_1(\Sigma),U(r))/U(r).
\end{equation}
対応は，ユニタリ表現 $\rho$ から平坦束 $E_\rho=\tilde\Sigma\times_\rho\C^r$ をとり，平坦接続の $(0,1)$ 部分を $\bar\partial_E$ とすることで与えられる．
\end{mythm}

（次数 $d\neq0$ の場合は，1点の周りのホロノミーを中心元 $e^{2\pi i d/r}$ とした射影ユニタリ表現に対応する．）
この定理は DUY 定理の Riemann 面版である：Riemann 面上では $\Lambda F=0$ と $F=0$ が同値だからである．
右辺は位相的・ゲージ理論的に定義された空間（複素構造によらない）であり，左辺は代数幾何的な空間（$\Sigma$ の複素構造に依存する）である．
同じ空間に二つの記述があることで，たとえば左辺から複素構造・射影多様体の構造が，右辺からシンプレクティック構造
（Atiyah--Bott：$\omega(\alpha,\beta)=-\int_\Sigma\Tr(\alpha\wedge\beta)$）が見え，$\mathcal{M}$ は Kähler 多様体となる．

\begin{note}
Atiyah--Bott の観点：Riemann 面上の接続の空間 $\mathcal{A}$ は無限次元のアフィン空間で，シンプレクティック形式
$\omega(\alpha,\beta)=-\int_\Sigma\Tr(\alpha\wedge\beta)$ をもつ．ゲージ群 $\mathcal{G}$ の作用の運動量写像は曲率 $F$ そのものであり，
$\mathcal{M}_{\rm flat}=F^{-1}(0)/\mathcal{G}$ はシンプレクティック商である．一方 $\mathcal{A}$ は正則構造 $\bar\partial_A$ の空間でもあり，
複素化ゲージ群 $\mathcal{G}_\C$ による商（安定な軌道に制限）が $\mathcal{N}$ である．
「シンプレクティック商 $=$ 安定点の複素商」は Kempf--Ness の定理の無限次元版であり，Narasimhan--Seshadri/DUY の概念的な理由である．
\end{note}

\subsection{Higgs束のモジュライ}

安定 Higgs 束 $(E,\Phi)$（階数 $r$，次数 $0$，行列式自明）の同型類全体 $\mathcal{M}_{\rm Higgs}$ は，複素次元
\begin{equation}
  \dim_\C\mathcal{M}_{\rm Higgs} = 2(r^2-1)(g-1)
\end{equation}
の準射影多様体である．安定束 $E$ に対し $\Phi\in H^0(\End_0 E\otimes K)$ は Serre 双対性により $H^1(\End_0E)^* = T^*_E\mathcal{N}$ に一致するので，
\begin{equation}
  T^*\mathcal{N}\subset\mathcal{M}_{\rm Higgs}
\end{equation}
は開稠密部分集合として含まれる．したがって $\mathcal{M}_{\rm Higgs}$ は（コンパクトな）束のモジュライの余接束を部分コンパクト化したもので，
正則シンプレクティック形式をもつ．

$\mathcal{M}_{\rm Higgs}$ は次の三つの顔をもつ（Hitchin）：
\begin{enumerate}
  \item \textbf{Dolbeault 描像}：安定 Higgs 束 $(E,\Phi)$ のモジュライ（代数幾何，$\Sigma$ の複素構造に依存）．
  \item \textbf{Hitchin 方程式の解}：$(A,\Phi)$ で \eqref{eq:ch14-hitchin} をみたすもの $/$ ユニタリゲージ群（ゲージ理論）．
  \item \textbf{de Rham／Betti 描像}：平坦 $\SL(r,\C)$ 接続のモジュライ（位相的）
  \begin{equation*}
    \mathcal{M}_{\rm dR}\cong\Hom(\pi_1(\Sigma),\SL(r,\C))/\!/\SL(r,\C).
  \end{equation*}
\end{enumerate}
(1) と (3) は複素多様体としては異なる（(3) はアフィン多様体，(1) は Hitchin fibration により固有なファイバーをもつ）が，
同じ実多様体の上の異なる複素構造であり，あわせて\textbf{超 Kähler 多様体}をなす．これは非可換 Hodge 対応の幾何学的内容である．

\begin{myex}{種数1・階数1}{ch14-higgs-moduli-ex}
$\Sigma=E_\tau$，$r=1$ では，Dolbeault 描像は $\mathrm{Jac}(E_\tau)\times H^0(K)\cong E_\tau\times\C$，
Betti 描像は $\Hom(\Z^2,\C^*)=\C^*\times\C^*$ である．両者は実多様体としては $T^2\times\R^2$ で同じだが，
複素多様体としては異なる（前者はコンパクトな因子 $E_\tau$ を含み，後者はアフィン）．
\end{myex}

\subsection{物理的自由度とモジュライ}

モジュライ空間は物理では\textbf{真空のモジュライ}，あるいは\textbf{BPS 解のモジュライ}として現れる．

\textbf{(1) 真空のモジュライと D 項・F 項}：
超対称ゲージ理論の古典的真空は，スカラー場 $\phi$ に対するポテンシャル $V = |\partial W/\partial\phi|^2 + \frac{g^2}{2}\sum_a(D^a)^2$ の零点，
すなわち F 項条件 $\partial W/\partial\phi=0$ と D 項条件 $D^a=\phi^\dagger T^a\phi=0$ の解をゲージ群 $G$ で割ったものである．
D 項条件は $G$ の作用の運動量写像の零点であり，Kempf--Ness の定理により
\begin{equation}
  \mathcal{M}_{\rm vac} = \{F\text{ 項}=0,\ D\text{ 項}=0\}/G \;\cong\; \{F\text{ 項}=0\}^{\rm (poly)stable}/G_\C
\end{equation}
と，複素化ゲージ群による（安定点の）代数的商として書ける．これは HYM 方程式と安定性の対応（DUY）の有限次元版であり，
「D 項 $\leftrightarrow$ 安定性」「F 項 $\leftrightarrow$ 正則性」という対応が，14.4.5節で注意したとおり一貫して現れる．
モジュライ空間の各点は異なる真空を表し，その接方向は質量をもたないスカラー場（モジュライ場）である．

\textbf{(2) BPS 方程式}：作用やエネルギーが位相不変量で下から押さえられ，等号が1階の方程式で実現されるとき，その方程式を BPS 方程式という．
\begin{center}
\small
\begin{tabular}{c|c|c}
BPS 方程式 & 次元 & モジュライ空間 \\ \hline
反自己双対方程式 $F=-*F$ & 4 & インスタントンのモジュライ（ADHM 構成） \\
Bogomolny 方程式 $F=*D\Phi$ & 3 & モノポールのモジュライ \\
Hitchin 方程式 \eqref{eq:ch14-hitchin} & 2 & $\mathcal{M}_{\rm Higgs}$ \\
渦糸方程式 & 2 & 対称積 $\mathrm{Sym}^k\Sigma$ \\
HYM 方程式 \eqref{eq:ch14-hym} & $2n$ & 安定束のモジュライ
\end{tabular}
\end{center}
これらはすべて4次元の反自己双対方程式の次元還元・一般化であり，共通して「正則性（$F^{0,2}=0$，$\bar\partial_A\Phi=0$）＋ 運動量写像の条件」
の形をもつ．その結果，モジュライ空間は代数幾何的な記述（安定な正則対象の商）をもつ．
たとえば渦糸方程式 $\bar\partial_A\phi=0$，$*iF = \frac{e^2}{2}(|\phi|^2-v^2)$ では，$\phi$ は直線束の正則切断であり，その零点（渦の位置）が解を決める：
$k$ 個の渦のモジュライ空間は，$\phi$ の零点集合，すなわち $\Sigma$ 上の次数 $k$ の有効因子の空間 $\mathrm{Sym}^k\Sigma$ に同型である（Bradlow，Taubes）．

\begin{physbox}[BPS 状態とモジュライ空間の量子化]
BPS 解の低エネルギーのダイナミクスは，モジュライ空間上の測地線運動（モジュライ空間近似，Manton）で記述される．
さらに超対称量子力学としてモジュライ空間を量子化すると，BPS 状態の縮退度はモジュライ空間のコホモロジー（あるいは指数）で数えられる．
たとえば D ブレーン上のゲージ理論では，ブレーンの束縛状態が安定層のモジュライ空間のコホモロジーで数えられ，
Donaldson--Thomas 不変量などの代数幾何的不変量として計算される（第17章）．
\end{physbox}

\exprob{ch14-pr-flatdim}
$G=SU(2)$ とし，$\Hom(\pi_1(\Sigma_g),G)/G$ の次元を「$2g$ 個の $G$ の元 $-$ 関係式（$G$ 値で $\dim G$ 個）$-$ 共役作用（$\dim G$）」として数え，
実次元が $6g-6$ であることを示せ．また $g=1$ の場合，$\pi_1(T^2)=\Z^2$ の $U(1)$ 表現の空間が双対トーラス $T^2$ であることを示せ．

\exprob{ch14-pr-vortex}
複素平面 $\C$ 上の $U(1)$ 渦糸方程式において，$k$ 個の渦の解が多項式 $\phi\sim\prod_{j=1}^k(z-z_j)$ の零点 $\{z_1,\dots,z_k\}$ で特徴づけられると仮定する．
このときモジュライ空間 $\mathrm{Sym}^k\C$ が $\C^k$ と同型であることを，多項式の係数を座標にとることで示せ．

\section*{問題の解答}
\addcontentsline{toc}{section}{問題の解答}

\exsol{ch14-pr-cocycle}
コサイクル条件 $g_{\alpha\beta}g_{\beta\gamma}=g_{\alpha\gamma}$ で $\alpha=\beta=\gamma$ とおくと $g_{\alpha\alpha}g_{\alpha\alpha}=g_{\alpha\alpha}$ であり，
$G$ の元 $g_{\alpha\alpha}(x)$ の逆元を右から掛けて $g_{\alpha\alpha}=e$．次に $\gamma=\alpha$ とおくと（$U_\alpha\cap U_\beta$ 上で）
$g_{\alpha\beta}g_{\beta\alpha}=g_{\alpha\alpha}=e$ なので $g_{\beta\alpha}=g_{\alpha\beta}^{-1}$．
（これらは定義 $s_\beta=s_\alpha g_{\alpha\beta}$ から直接も分かる：$s_\alpha=s_\alpha g_{\alpha\alpha}$ と作用の自由性から $g_{\alpha\alpha}=e$．）

局所切断を $s'_\alpha=s_\alpha h_\alpha$（$h_\alpha:U_\alpha\to G$）に取り替えると
\begin{equation*}
  s'_\beta = s_\beta h_\beta = s_\alpha g_{\alpha\beta}h_\beta = s'_\alpha\,h_\alpha^{-1}g_{\alpha\beta}h_\beta
\end{equation*}
より新しい遷移関数は $g'_{\alpha\beta}=h_\alpha^{-1}g_{\alpha\beta}h_\beta$ である．三重の重なりで
\begin{equation*}
  g'_{\alpha\beta}g'_{\beta\gamma} = h_\alpha^{-1}g_{\alpha\beta}h_\beta h_\beta^{-1}g_{\beta\gamma}h_\gamma = h_\alpha^{-1}g_{\alpha\gamma}h_\gamma = g'_{\alpha\gamma}
\end{equation*}
となり，$g'$ もコサイクル条件をみたす（したがって前半から $g'_{\alpha\alpha}=e$，$g'_{\beta\alpha}=g'^{-1}_{\alpha\beta}$ も成り立つ．直接にも $g'_{\alpha\alpha}=h_\alpha^{-1}h_\alpha=e$）．

\exsol{ch14-pr-od}
$U_0$ の座標を $z=Z_1/Z_0$，$U_1$ の座標を $w=Z_0/Z_1$ とする．本文より $\OO(d)$ の遷移関数は $g_{01}=(Z_1/Z_0)^d=z^d$ であり，
正則切断 $s$ の成分 $s_0(z)$（$\C$ 上正則），$s_1(w)$（$\C$ 上正則）は \eqref{eq:ch14-section-transform} により
\begin{equation*}
  s_0(z) = z^d\,s_1(1/z)\qquad(z\neq0)
\end{equation*}
で結ばれる．冪級数展開 $s_0(z)=\sum_{k\geq0}a_kz^k$，$s_1(w)=\sum_{m\geq0}b_mw^m$ を代入すると，$\C^*$ 上で
\begin{equation*}
  \sum_{k\geq0}a_kz^k = \sum_{m\geq0}b_mz^{d-m}
\end{equation*}
であり，Laurent 展開の一意性から $a_k=b_{d-k}$ かつ，$z^k$ の係数は $k\geq0$ と $d-k\geq0$ の両方をみたすときのみ $0$ でありうる．
\begin{itemize}
  \item $d<0$ のとき：$0\leq k\leq d$ をみたす $k$ はないので $a_k=b_m=0$，すなわち $s=0$．よって $H^0(\PP^1,\OO(d))=0$．
  \item $d\geq0$ のとき：$s_0$ は次数 $d$ 以下の多項式 $\sum_{k=0}^da_kz^k$ であり，$s_1(w)=\sum_{k=0}^d a_kw^{d-k}$ で決まる．
  $P(Z_0,Z_1):=Z_0^d\,s_0(Z_1/Z_0)=\sum_{k=0}^da_kZ_0^{d-k}Z_1^k$ とおけば $P$ は $d$ 次同次多項式で，$s_0=P/Z_0^d$，$s_1=P/Z_1^d$ である
  （実際 $P/Z_0^d = (Z_1/Z_0)^d\cdot P/Z_1^d$ は $s_0=z^ds_1$ と一致する）．
  逆に任意の $d$ 次同次多項式 $P$ から $s_\alpha=P/Z_\alpha^d$ により正則切断が得られる．
\end{itemize}
対応 $s\leftrightarrow P$ は線形な全単射であり，$d$ 次同次多項式の空間の基底は $Z_0^{d-k}Z_1^k$（$k=0,\dots,d$）なので
$\dim H^0(\PP^1,\OO(d))=d+1$ である．（本文の公式 $\binom{n+d}{d}$ の $n=1$ の場合．）

\exsol{ch14-pr-mobius}
\textbf{切断の零点}：$s$ を $M$ の切断とし，局所成分を $s_1:U_1\to\R$，$s_2:U_2\to\R$ とする．$s_1=g_{12}s_2$ より
$I_+$ 上で $s_1=s_2$，$I_-$ 上で $s_1=-s_2$ である．$s$ が零点をもたないと仮定すると，弧 $U_1,U_2$ は連結なので
中間値の定理により $s_1,s_2$ はそれぞれ一定の符号をもつ．$I_+$ 上の関係から $s_1$ と $s_2$ の符号は等しく，$I_-$ 上の関係から逆である．矛盾．
よって任意の切断は零点をもつ（図\ref{fig:ch14-mobius}）．とくに零点のない切断（＝大域的な枠）がないので $M$ は自明でない．

\textbf{$M\oplus M$ の自明性}：$M\oplus M$ の遷移関数は $g_{12}\oplus g_{12}=\pm\mathbf{1}_2$ である．一方，$\R^2=\C$ と同一視すると，
複素直線束 $M_\C:=M\otimes_\R\C$（遷移関数 $\pm1\in\C^*$）は実階数2の束として遷移関数 $\pm\mathbf{1}_2$ をもつので，$M\oplus M\cong M_\C$ である．
$M_\C$ が自明であることを示す．$S^1=\R/\Z$ の座標 $t$ を用い，$U_1=(-\epsilon,\tfrac12+\epsilon)$，$U_2=(\tfrac12-\epsilon,1+\epsilon)$（$\bmod 1$），
$I_+=(\tfrac12-\epsilon,\tfrac12+\epsilon)$，$I_-=(-\epsilon,\epsilon)$ とする．滑らかな関数 $\psi:U_1\to[0,1]$ を
$t<\epsilon$ で $\psi=0$，$t>\tfrac12-\epsilon$ で $\psi=1$ となるようにとり，
\begin{equation*}
  h_1(t) = -e^{i\pi\psi(t)}\in U(1)\subset\C^*\quad(U_1\text{ 上}),\qquad h_2=1\quad(U_2\text{ 上})
\end{equation*}
とおく．$I_-$ 上で $h_1=-1$，$I_+$ 上で $h_1=-e^{i\pi}=1$ であるから，枠の取り替え後の遷移関数（問\ref{ch14-pr-cocycle}）は
\begin{equation*}
  g'_{12}=h_1^{-1}g_{12}h_2 = \begin{cases}1^{-1}\cdot1\cdot1=1 & (I_+)\\ (-1)^{-1}\cdot(-1)\cdot1=1 & (I_-)\end{cases}
\end{equation*}
となり，恒等的に $1$ である．よって $M_\C$ は自明な複素直線束 $S^1\times\C$ と同型であり，実束として $M\oplus M\cong S^1\times\R^2$．
（実の言葉では，$-\mathbf{1}_2$ は角度 $\pi$ の回転であり，$\mathbf{1}_2$ と $\GL^+(2,\R)$ の中で回転の道 $e^{i\pi\psi}$ によって結べることが本質である．
$\GL(1,\R)=\R^*$ は連結でないので $M$ 自身ではこれができない．）

\exsol{ch14-pr-curv}
$A'=g^{-1}Ag+g^{-1}dg$ とおく．$gg^{-1}=1$ を微分して $d(g^{-1})=-g^{-1}dg\,g^{-1}$ である．1形式 $A$ について $d(Ag)=dA\,g-A\wedge dg$ に注意して
\begin{align*}
  dA' &= d(g^{-1})\wedge Ag + g^{-1}dA\,g - g^{-1}A\wedge dg + d(g^{-1})\wedge dg\\
      &= -g^{-1}dg\,g^{-1}\wedge Ag + g^{-1}dA\,g - g^{-1}A\wedge dg - g^{-1}dg\,g^{-1}\wedge dg ,\\
  A'\wedge A' &= (g^{-1}Ag+g^{-1}dg)\wedge(g^{-1}Ag+g^{-1}dg)\\
      &= g^{-1}A\wedge A\,g + g^{-1}A\wedge dg + g^{-1}dg\,g^{-1}\wedge Ag + g^{-1}dg\,g^{-1}\wedge dg .
\end{align*}
和をとると，$g^{-1}dg\,g^{-1}\wedge Ag$，$g^{-1}A\wedge dg$，$g^{-1}dg\,g^{-1}\wedge dg$ の項がそれぞれ打ち消し合い，
\begin{equation*}
  F' = dA'+A'\wedge A' = g^{-1}dA\,g + g^{-1}A\wedge A\,g = g^{-1}(dA+A\wedge A)g = g^{-1}Fg
\end{equation*}
を得る．（$U(1)$ では $g^{-1}Fg=F$ であり，場の強さはゲージ不変である．）

\exsol{ch14-pr-mono}
\textbf{滑らかさ}：$A_N$ の係数 $1-\cos\theta$ は北極 $\theta=0$ で $\theta^2/2+O(\theta^4)$ のように消える．北極の近くの局所座標 $x=\theta\cos\varphi$，$y=\theta\sin\varphi$ では
$\theta^2d\varphi=x\,dy-y\,dx$ であり，$(1-\cos\theta)/\theta^2$ は $\theta^2=x^2+y^2$ の滑らかな関数なので，$A_N$ は北極を含む $U_N$ 上で滑らかである．
同様に $1+\cos\theta$ は南極 $\theta=\pi$ で2位で消えるので，$A_S$ は $U_S$ 上で滑らかである．
また $A_N,A_S$ は純虚数値（$\mathfrak{u}(1)=i\R$ 値）であり，ユニタリ接続を表す．

\textbf{貼り合わせ}：$U(1)$ は可換なので \eqref{eq:ch14-A-transform}（$\alpha=N$，$\beta=S$，$g=g_{NS}=e^{in\varphi}$）は
$A_S=A_N+g_{NS}^{-1}dg_{NS}=A_N+in\,d\varphi$ となる．実際，重なり $U_N\cap U_S=\{0<\theta<\pi\}$ で
\begin{equation*}
  A_S - A_N = \frac{in}{2}(1+\cos\theta)\,d\varphi + \frac{in}{2}(1-\cos\theta)\,d\varphi = in\,d\varphi = e^{-in\varphi}\,d(e^{in\varphi})
\end{equation*}
であり，確かに成り立つ．$e^{in\varphi}$ が赤道上で1価であるために $n\in\Z$ が必要である（例\ref{ex:ch14-dirac}）．
（$\alpha=S,\beta=N$ と読むと $A_N=A_S+g_{SN}^{-1}dg_{SN}$，$g_{SN}=g_{NS}^{-1}=e^{-in\varphi}$ となり，同じことである．
$\alpha,\beta$ の役割を取り違えて $g=e^{in\varphi}$ を当てはめると $n$ の符号が逆転するので注意．）

\textbf{曲率と積分}：$U(1)$ では $A\wedge A=0$ なので $F=dA$ であり，
\begin{equation*}
  F|_{U_N}=dA_N = -\frac{in}{2}\sin\theta\,d\theta\wedge d\varphi,\qquad
  F|_{U_S}=dA_S = \frac{in}{2}(-\sin\theta)\,d\theta\wedge d\varphi = -\frac{in}{2}\sin\theta\,d\theta\wedge d\varphi .
\end{equation*}
両者は重なりで一致し（$d(in\,d\varphi)=0$），$F$ は $S^2$ 上の大域的な2形式である．$\int_{S^2}\sin\theta\,d\theta\wedge d\varphi=4\pi$ より
\begin{equation*}
  \int_{S^2}F = -\frac{in}{2}\cdot4\pi = -2\pi i n,\qquad \frac{i}{2\pi}\int_{S^2}F = \frac{i}{2\pi}(-2\pi in) = n .
\end{equation*}
Stokes の定理でも確かめられる：北半球 $H_N$ と南半球 $H_S$ に分け，赤道 $\gamma$ を $H_N$ の境界の向き（$\varphi$ 増加方向）にとると
$\int_{S^2}F=\oint_\gamma A_N-\oint_\gamma A_S=-\oint_\gamma in\,d\varphi=-2\pi in$．積分は遷移関数の巻き数だけで決まる．

\textbf{物理的な意味と $\OO(-1)$ との整合}：\eqref{eq:ch14-convention} より $A^{\rm phys}=\frac{i}{g}A$ なので
$A^{\rm phys}_N=\frac{n}{2g}(1-\cos\theta)\,d\varphi$ であり，これは磁荷 $g_m=n/(2g)$ の Dirac モノポールの標準的なポテンシャルである．
全磁束は $\int F^{\rm phys}=4\pi g_m=2\pi n/g$，すなわち Dirac の量子化条件 $g\Phi_B=2\pi n$ であり，
$\frac{i}{2\pi}\int F=\frac{g}{2\pi}\int F^{\rm phys}=n$ はまさに磁荷（Dirac 単位）である．

例\ref{ex:ch14-o-1} の $\OO(-1)$ を見ると，$z=\tan(\theta/2)e^{i\varphi}$ で $U_0=\{Z_0\neq0\}$ を $U_N$ に（$z=0$ が北極），$U_1$ を $U_S$ に対応させる．
正則枠 $e_0,e_1$（$e_1=e_0g_{01}$，$g_{01}=Z_0/Z_1=1/z$）を正規化したユニタリ枠 $u_\alpha=e_\alpha/|e_\alpha|$ の間の遷移関数は，$|e_0|=\sqrt{1+|z|^2}$，$|e_1|=\sqrt{1+|z|^{-2}}=|e_0|/|z|$ より
\begin{equation*}
  u_1 = u_0\cdot\frac{1}{z}\cdot\frac{|e_0|}{|e_1|} = u_0\cdot\frac{|z|}{z} = u_0\,e^{-i\varphi}
\end{equation*}
であり，$g_{NS}=e^{-i\varphi}$，すなわち $n=-1$ である．$z=\tan(\theta/2)e^{i\varphi}$ の向き $dx\wedge dy$ は $d\theta\wedge d\varphi$ と同じ向きなので，
上の結果 $\frac{i}{2\pi}\int F=n=-1$ は例\ref{ex:ch14-o-1} の計算 $\frac{i}{2\pi}\int_{\PP^1}F=-1$ と一致する．
こうして「$\frac{i}{2\pi}\int F$ ＝ 第1 Chern 数 ＝ 磁荷 $n$」，$\OO(-1)$ は磁荷 $-1$（$\OO(1)$ が磁荷 $+1$）という本文の規約が整合的に確かめられた．

\exsol{ch14-pr-chern}
正則枠 $e$ に関して $H=h(e,e)=e^{-\varphi}$（$1\times1$ 行列）である．Chern 接続の公式 \eqref{eq:ch14-chern} より
\begin{equation*}
  A = H^{-1}\partial H = e^{\varphi}\,\partial(e^{-\varphi}) = e^{\varphi}\cdot(-e^{-\varphi}\partial\varphi) = -\partial\varphi,
\end{equation*}
\begin{equation*}
  F = \bar\partial A = -\bar\partial\partial\varphi = \partial\bar\partial\varphi
\end{equation*}
（$\partial\bar\partial=-\bar\partial\partial$ を用いた）．直線束なので $A\wedge A=0$ であり，$F=dA=\partial A+\bar\partial A$ で $\partial A=-\partial^2\varphi=0$ からも同じ結果が出る．

念のため定義から直接確かめる．$A=-\partial\varphi$ は $(1,0)$ 型なので $D^{0,1}e=0=\bar\partial_Ee$ であり，また $D e = A\,e$ より
\begin{equation*}
  h(De,e)+h(e,De) = (A+\bar A)H = -(\partial\varphi+\bar\partial\varphi)e^{-\varphi} = -d\varphi\,e^{-\varphi} = dH
\end{equation*}
（$\varphi$ は実数値なので $\overline{\partial\varphi}=\bar\partial\varphi$）となって両立条件もみたす．一意性からこれが Chern 接続である．

例：$\OO(-1)$ では $H=1+|z|^2$，$\varphi=-\log(1+|z|^2)$ であり，$F=\partial\bar\partial\varphi=-\partial\bar\partial\log(1+|z|^2)=\bar\partial\partial\log(1+|z|^2)$ は例\ref{ex:ch14-o-1} と一致する．

\exsol{ch14-pr-stab}
\textbf{直線束}：$E=L$ を直線束とすると $\rank E=1$ であり，安定性の定義で検査すべき部分層 $\sheaf{F}$ は $0<\rank\sheaf{F}<1$ をみたさねばならないが，
そのような整数 $\rank\sheaf{F}$ は存在しない．したがって条件は空であり，$L$ は自動的に安定である．
（直観的には，直線束には「より小さい傾きをもつべき真の部分」がない．ゲージ理論側では，Riemann 面上の直線束は定数曲率の Hermitian 計量，
すなわち HYM 接続を常にもつことに対応する．）

\textbf{$\OO(1)\oplus\OO(-1)$}：$E=\OO(1)\oplus\OO(-1)$ は $\rank E=2$，$\deg E=1+(-1)=0$ なので $\mu(E)=0$．
第1成分 $\sheaf{F}=\OO(1)\subset E$ は階数1の部分束（とくに連接部分層）で $\mu(\sheaf{F})=1$．
$\mu(\sheaf{F})=1>0=\mu(E)$ なので $E$ は安定でなく，半安定ですらない．すなわち不安定である（例\ref{ex:ch14-p1-stab} の $a=1,b=-1$ の場合）．
DUY 定理により，$E$ は（$\OO(1)$ と $\OO(-1)$ の直和として傾きの異なる成分をもち）HYM 接続をもたない．

\exsol{ch14-pr-hitchin-dim}
$\deg K=2g-2$ より $\deg K^k=k(2g-2)$ である．$k\geq2$，$g\geq2$ なら
\begin{equation*}
  \deg K^k - (2g-2) = (k-1)(2g-2) > 0
\end{equation*}
なので $H^1(\Sigma,K^k)=0$．Riemann--Roch より
\begin{equation*}
  \dim H^0(\Sigma,K^k) = \deg K^k+1-g = 2k(g-1)-(g-1) = (2k-1)(g-1).
\end{equation*}
（$H^1(L)=0$ は Serre 双対性 $H^1(L)\cong H^0(K\otimes L^{-1})^*$ と，負の次数の直線束が正則切断をもたないことから従う．
$k=1$ では $H^1(K)\cong H^0(\OO)^*=\C$ なので公式は $\dim H^0(K)=g$ となり，$(2k-1)(g-1)$ とは異なる．）

$\SL(r)$ Hitchin 系の底空間は $a_1=\Tr\Phi=0$ として $\mathcal{B}=\bigoplus_{k=2}^rH^0(\Sigma,K^k)$ であるから
\begin{equation*}
  \dim\mathcal{B} = \sum_{k=2}^r(2k-1)(g-1) = (g-1)\Bigl(\sum_{k=1}^r(2k-1)-1\Bigr) = (g-1)(r^2-1).
\end{equation*}
これは $\dim_\C\mathcal{M}_{\rm Higgs}=2(r^2-1)(g-1)$ のちょうど半分であり，Hitchin fibration が完全可積分系（Lagrange ファイブレーション）であることと整合する．
$r=2$ では $3(g-1)=3g-3$ で本文の値に一致する．

\exsol{ch14-pr-flatdim}
\textbf{次元の数え上げ}：$\Hom(\pi_1(\Sigma_g),G)$ は $G^{2g}$ の中で関係式
\begin{equation*}
  \mu(A_1,B_1,\dots,A_g,B_g) := \prod_{i=1}^gA_iB_iA_i^{-1}B_i^{-1} = 1
\end{equation*}
で切り出される．$G^{2g}$ の実次元は $2g\dim G$，関係式は $G$ に値をとるので $\dim G$ 個の実条件を課す．
既約表現（中心化群が $G$ の中心 $\{\pm1\}$ のみ，$g\geq2$ で一般の点）では $\mu$ の微分が全射で，$\mu^{-1}(1)$ はその近くで実次元
$2g\dim G-\dim G$ の多様体となる．さらに $G$ の共役作用は中心が自明に作用するのを除いて局所的に自由なので，商をとると $\dim G$ だけ次元が下がる．
$G=SU(2)$ では $\dim G=3$ なので
\begin{equation*}
  \dim_\R\Hom(\pi_1(\Sigma_g),SU(2))/SU(2) = 2g\cdot3-3-3 = 6g-6
\end{equation*}
である（$g\geq2$）．複素次元に直すと $3g-3=(r^2-1)(g-1)|_{r=2}$ で，本文や問\ref{ch14-pr-hitchin-dim} の $\SL(2)$ の場合と一致する．

\textbf{$g=1$，$G=U(1)$}：$\Sigma=T^2=\R^2/\Lambda$，$\Lambda=\Z\lambda_1+\Z\lambda_2$ とすると $\pi_1(T^2)=\Lambda\cong\Z^2$ は可換で，
関係式 $ABA^{-1}B^{-1}=1$ は $U(1)$ が可換なので自動的にみたされる．また $U(1)$ の共役作用は自明である．よって
\begin{equation*}
  \Hom(\pi_1(T^2),U(1))/U(1) = \Hom(\Lambda,U(1)) \cong U(1)\times U(1),\qquad \rho\mapsto(\rho(\lambda_1),\rho(\lambda_2))
\end{equation*}
（$\Lambda$ は自由加群なので，生成元の像を任意に指定できる）．これは2次元トーラスである．
より内在的には，$\theta\in(\R^2)^*$ に指標 $\rho_\theta(\lambda)=e^{2\pi i\langle\theta,\lambda\rangle}$ を対応させる写像 $(\R^2)^*\to\Hom(\Lambda,U(1))$ は
全射準同型（$\langle\theta,\lambda_j\rangle$ を任意に選べる）で，その核は双対格子 $\Lambda^*=\{\theta:\langle\theta,\lambda\rangle\in\Z\ (\forall\lambda\in\Lambda)\}$ である．したがって
\begin{equation*}
  \mathcal{M}_{\rm flat}(T^2,U(1)) \cong (\R^2)^*/\Lambda^* = \widehat{T^2}
\end{equation*}
は双対トーラスである．$\rho_\theta$ は定数係数の平坦接続 $A=-2\pi i\langle\theta,dx\rangle$ のホロノミー $\exp(-\oint_\lambda A)=e^{2\pi i\langle\theta,\lambda\rangle}$ として実現され，
$\theta\mapsto\theta+\Lambda^*$ は大域的なゲージ変換 $e^{2\pi i\langle\theta',x\rangle}$（$\theta'\in\Lambda^*$，$T^2$ 上1価）に対応する．
（$T^2$ の場合は表現がすべて可約なので，上の $6g-6$ の数え方は適用できない．）

\exsol{ch14-pr-vortex}
$\mathrm{Sym}^k\C=\C^k/S_k$（$S_k$ は座標の置換で作用）とし，写像
\begin{equation*}
  \sigma:\C^k\to\C^k,\qquad (z_1,\dots,z_k)\mapsto(a_1,\dots,a_k),\qquad \prod_{j=1}^k(z-z_j) = z^k+a_1z^{k-1}+\cdots+a_k
\end{equation*}
を考える．$a_m=(-1)^me_m(z_1,\dots,z_k)$（$e_m$ は $m$ 次基本対称式）は多項式で $S_k$ 不変なので，$\sigma$ は連続（正則）写像
$\bar\sigma:\mathrm{Sym}^k\C\to\C^k$ を誘導する．
\begin{itemize}
  \item \textbf{全射}：任意の $(a_1,\dots,a_k)$ に対し，代数学の基本定理によりモニック多項式 $z^k+a_1z^{k-1}+\cdots+a_k$ は $\prod_j(z-z_j)$ と1次式に分解する．
  \item \textbf{単射}：多項式環 $\C[z]$ は一意分解整域なので，根 $\{z_1,\dots,z_k\}$ は重複度込みで（順序を除いて）一意である．すなわち $\sigma(z)=\sigma(z')$ なら $z'$ は $z$ の置換である．
  \item \textbf{同相}：$\sigma$ は固有写像である（根は $|z_j|\leq1+\max_m|a_m|$ で押さえられるので，有界集合の逆像は有界）．
  したがって $\bar\sigma$ は固有な連続全単射で，局所コンパクト Hausdorff 空間の間の写像として閉写像，よって同相である．
\end{itemize}
さらに代数的には，対称式の基本定理 $\C[z_1,\dots,z_k]^{S_k}=\C[e_1,\dots,e_k]$（$e_m$ は代数的に独立）により
\begin{equation*}
  \mathrm{Sym}^k\C = \C^k/S_k = \Spec\C[z_1,\dots,z_k]^{S_k} = \Spec\C[a_1,\dots,a_k]\cong\A^k
\end{equation*}
であり，$\bar\sigma$ は代数多様体（したがって複素多様体）としての同型である．こうして $k$ 個の渦のモジュライ空間は，
多項式 $\phi$ の係数 $(a_1,\dots,a_k)$ を大域座標とする $\C^k$ である．
渦が重なる点（$z_i=z_j$）は $\C^k/S_k$ の商としては特異に見えるが，係数座標では何ら特異でないことに注意する（これは1次元の底空間に特有の現象である）．
